% =============================================================================
% Appendix C: self-contained continuity proof (inlined June 2026).
% Source: theory/continuity appendix/continuity_appendix.tex (standalone).
% Transformations applied when inlining:
%   - standalone preamble / title / abstract / thebibliography / \end{document} stripped;
%   - section levels demoted by one (\section->\subsection, \subsection->\subsubsection);
%   - citation keys unified to references.bib (codina2001->codina2001stabilized,
%     brennerscott->brenner2008mathematical; ernguermond added to references.bib);
%   - the standalone "Amendments to the original statements" section was dropped
%     (its main-text amendments are now applied in the body of the article);
%   - macros and the assumption environment are declared in the article preamble.
%
% Restructured 2026-07-28 to run parallel to Appendix D (osgs_appendix.tex): setting ->
% standing assumptions -> technical results (parameters, weighted inverse estimates, jump,
% global identities, triple norm) -> stability -> continuity -> convergence.  The duplicated
% "Proof of coercivity" subsubsection (label sec:proof-stability, display eq:coerccollect-lemma)
% was removed; its content lives once, in full, in the Stability subsection.
%
% [known-fragility]  This file is \input by BOTH article.tex (v1) and article_v2.tex (v2).
% It may therefore cite only labels defined in BOTH.  In particular do NOT cite
% sec:CommonSetting, eq:TripleNormASGS, eq:GalerkinCoercivityIdentity, eq:CollectedCoercivity,
% app:osgs, sec:StabilityOSGS or any oa:* / H:design.. label: all are v2-only and silently
% become undefined references in the v1 build.
%
% NOTATION SET (consolidated into Sec. 3.2 on 2026-07-29; cite these two, restate nothing):
%   eq:TauInnerProduct  (Eq. 3.16) -- the ONE definitional display: line 1 the tau-matrix
%                                     product and its norm, line 2 the broken (.,.)_h and
%                                     ||.||_h, plus a prose comment right after it covering
%                                     the single-parameter weighting (.,.)_{tau_i},
%                                     ||.||_{tau_i}.  CITE THIS for the notation set.
%   eq:L2InnerProduct   (Eq. 3.17) -- ONLY the approximation (.,.)_tau ~~ (.,.)_h, i.e. the
%                                     implemented plain-L2 projection.  It defines nothing;
%                                     cite it only where the implemented-vs-analyzed
%                                     distinction is the point (see oa:rem:analyzed).
% Two Sec. 5 displays that duplicated these were RETIRED on 2026-07-29: eq:TauWeightedProduct
% (it re-defined, with \coloneqq, the product already at eq:TauInnerProduct) and
% eq:BrokenProduct (its RHS is literally the RHS of eq:L2InnerProduct).  Do not reintroduce
% either; both mains now carry the whole set in Sec. 3.2.
%
% Labels this file cites that are GUARANTEED to exist in both builds (keep in sync):
%   H:mesh, H:data, H:projector, H:porosity, H:advection, H:spaces, H:regularity,
%   H:coercivity, H:jump; eq:PiKorn, eq:L2InnerProduct, eq:TauInnerProduct,
%   eq:Tau1Final, eq:Tau2Final, eq:TauNavierStokes, eq:UpperBoundOnEpsilon,
%   eq:JumpCondition, eq:SigmaAlpha, eq:resolved, eq:conditions_on_num_param,
%   eq:InterpolationError, eq:ConvergenceResult, sec:StabilityASGS,
%   sec:TheContinuousProblem, sec:ViscousProjector, sec:NumericalExamplesStationaryCase3D,
%   sec:Conclusions,
%   lemma:Stability, lemma:Continuity, th:Convergence,
%   lem:projfamily, lem:projinstances.
%   (The rest of the relocated projector labels -- eq:PiOrthogonal, eq:projmonotone,
%   eq:devsymidentity, eq:crossibp, eq:projconstants -- are
%   NOT cited from here; they live in sec:ViscousProjector of both mains.  Add a name
%   to the list above only when this file actually starts citing it.)
%   Labels this file DEFINES that BOTH mains cite (keep them defined here, and keep the
%   wording usable by both builds -- no v2-only notation):
%     rem:porosityresolution -- cited from item H:porosity of each main's assumption list
%       (2026-07-30: the extended discussion of (A4) was moved out of that list into this
%       remark; it must therefore stay free of any reference to H:patch / chi_1, which is
%       v2-only, and states the surrogate caveat with locally introduced symbols instead).
%   (H:projector / eq:PiKorn were added to BOTH mains on 2026-07-29, immediately after
%   the H:data item, so that every existing \crefrange / dash range over the standing
%   assumptions keeps bracketing them without edits.)
%
% 2026-07-29 (RELOCATION): the subsubsection "Admissible viscous projectors"
% (\label{sec:projectors}: lem:projfamily, lem:projinstances and their
% displays) was MOVED OUT of this appendix into a new main-text subsection
% \label{sec:ViscousProjector} ("The viscous projector", Sec. 2.1), present VERBATIM in
% BOTH mains -- which is what keeps the labels above resolvable in both builds.  The
% display eq:PiKorn moved with it, out of the (A3) item and into that subsection, so (A3)
% now states the bare conditions and cites eq:PiOrthogonal / eq:PiKorn.  Do NOT
% reintroduce the block here: a second definition of those labels makes them
% multiply-defined in both builds.  What stays here is only the appendix-local
% bookkeeping (the setting sentence in Sec. C.1 and rem:hypotheses).
% =============================================================================
\section{Convergence analysis of the ASGS method}\label{app:Continuity}

This appendix establishes, for the linearized ASGS method, the stability
(\cref{lem:coercivity}), continuity (\cref{lem:continuity}) and convergence
(\cref{thm:convergence}) results stated as
\cref{lemma:Stability,lemma:Continuity,th:Convergence} in \cref{sec:StabilityASGS}. It follows
the algebraic{\hyp}subgrid{\hyp}scale template of~\cite{codina2001stabilized}, transported to the
porosity{\hyp}weighted operator, the Darcy reaction term and the artificial{\hyp}compressibility
perturbation. The proof is developed in full; only the material shared with the main text is
not re{\hyp}derived. Specifically, the continuous problem, the Galerkin bilinear form and its
coercivity identity, the stabilization parameters and the standing assumptions are stated once
in \cref{sec:StabilityASGS} and are here only recalled and referenced, so that this appendix
develops the ingredients specific to the ASGS route---the porosity{\hyp}weighted inverse
estimates, the coercivity of the stabilized form obtained by testing with the solution itself,
the eighteen{\hyp}term continuity estimate with its elementwise integrations by parts and jump
machinery, and the convergence theorem that follows from the two.

\subsection{The linearized problem and the ASGS method}\label{sec:setting}

Let $\Omega\subset\mathbb{R}^d$, $d\in\{2,3\}$, be a bounded, connected polyhedral domain with
boundary $\Gamma$. As in the analysis section of the main text, we consider the ASGS
method applied to the linearized porous Navier--Stokes problem with homogeneous Dirichlet
boundary conditions on the whole of $\Gamma$, uniform kinematic viscosity $\nu>0$ and a
uniform, isotropic resistance $\boldsymbol{\sigma}=\sigma\mathbb{I}$, $\sigma\ge 0$. Throughout,
$\ba$ is the given advection field, $\alpha:\Omega\to(0,1]$ the fluid volume fraction,
$\varepsilon\ge 0$ the penalty parameter, and $\ViscProj$ any viscous projector admissible in
the sense of \cref{sec:ViscousProjector}, as \cref{H:projector} requires; the
deviatoric--symmetric default $\ViscProj=\DSPi$ of \cref{sec:TheContinuousProblem} is one
instance, and nothing below is specific to it. \Cref{rem:hypotheses} records where each half of
\cref{H:projector} enters the proofs that follow.

The linearized differential operator $\mathcal{L}_{\ba}$ of \cref{sec:StabilityASGS}, acting on a
pair $U=[\bu;p]$, and its formal adjoint, acting on $V=[\bv;q]$, have momentum and mass
components
\begin{equation}\label{eq:StrongOperator}
  \mathcal{L}_{\ba}U =
  \begin{bmatrix}
    \alpha\X(U)-G(\bu)\\[4pt]
    \varepsilon p + \nabla\cdot(\alpha\bu)
  \end{bmatrix},
  \qquad
  \mathcal{L}_{\ba}^{*}V =
  \begin{bmatrix}
    -\alpha\X(V)-G(\bv)\\[4pt]
    \varepsilon q - \nabla\cdot(\alpha\bv)
  \end{bmatrix},
\end{equation}
where, on any pair $[\bv;q]$,
\begin{equation}\label{eq:XG}
  \X(V) = \ba\cdot\nabla\bv + \nabla q,
  \qquad
  G(\bv) \coloneqq 2\nu\,\nabla\cdot\bigl(\alpha\ViscProj\nabla\bv\bigr) - \sigma\bv ,
\end{equation}
the same formulas giving $\X(U)$ and $G(\bu)$ in the first column. Here $\X$ is the
first{\hyp}order momentum operator of \cref{sec:StabilityASGS} and $G$ collects the viscous and
reactive terms; keeping the two apart is what later turns the $\tau_1$ product into a difference
of squares \eqref{eq:stabdiag1}. The adjoint is written under $\nabla\cdot(\alpha\ba)=0$
(\cref{H:advection}) and with the simplification of the main text---the removal of the term
proportional to $\nabla\cdot(\alpha\ba)$, which vanishes anyway in the present setting. Passing to
it then reverses the sign of the first{\hyp}order terms---$\alpha\X$ in the momentum row,
$\nabla\cdot(\alpha\,\cdot)$ in the mass row---and leaves $G$ and the $\varepsilon$ terms
unchanged.

Let $\Th$ be a finite element partition of $\overline{\Omega}$ of elements $K$ with
diameters $h_K$, and let
$\Xh = \mathcal{V}_{h0}\times\mathcal{Q}_{h0}$, with
$\mathcal{V}_{h0}\subset H_0^1(\Omega)^d$ and
$\mathcal{Q}_{h0}\subset H^1(\Omega)\cap\mathcal{Q}_0$, be the continuous finite element
spaces of polynomial degrees $k_u\ge 1$ and $k_p\ge 1$ used for the velocity and the
pressure ($\mathcal{Q}_0$ being $L^2(\Omega)$, or its zero-mean subspace when
$\varepsilon=0$). The $L^2(\omega)$ inner product is written
$(\cdot,\cdot)_\omega$, with $\omega$ omitted when $\omega=\Omega$; elementwise defined
quantities are assembled in the broken and $\boldsymbol{\tau}${\hyp}weighted products and norms
\cref{eq:TauInnerProduct} of the main text, of whose two spellings the proofs below use the
expanded ones, $(f,\tau_i\,g)_h$ and $\bigl\lVert\tau_i^{1/2}f\bigr\rVert_h$, throughout.

Restricted to this setting, the stabilized bilinear form of the main text,
$\BASGS (\ba; U_h,V_h)=B(\ba; U_h,V_h)-\sum_K \langle \mathcal{L}_{\ba}^{*} V_h,
\boldsymbol{\tau}_K \mathcal{L}_{\ba} U_h\rangle_K$, reads, for
$U_h=[\buh;p_h]$ and $V_h=[\bvh;q_h]$ in $\Xh$,
\begin{equation}\label{eq:Bstab}
\begin{aligned}
  \BASGS (\ba; U_h,V_h)
  ={}& 2\nu\bigl(\ViscProj\nabla\bvh,\alpha\ViscProj\nabla\buh\bigr)
     + \bigl(\bvh,\alpha\X(U_h)\bigr)
     + \sigma\bigl(\bvh,\buh\bigr) \\
   & + \varepsilon\bigl(q_h,p_h\bigr)
     + \bigl(q_h,\nabla\cdot(\alpha\buh)\bigr) \\
   & + \bigl(\alpha\X(V_h)+G(\bvh),\,
       \tau_1\bigl[\alpha\X(U_h)-G(\buh)\bigr]\bigr)_h \\
   & + \bigl(\nabla\cdot(\alpha\bvh)-\varepsilon q_h,\,
       \tau_2\bigl[\varepsilon p_h+\nabla\cdot(\alpha\buh)\bigr]\bigr)_h ,
\end{aligned}
\end{equation}
together with
$\LASGS(\ba;V_h)=(\bvh,\bff)
 +\bigl(\alpha\X(V_h)+G(\bvh),\tau_1\bff\bigr)_h$.
For the exact solution $U$ of the linearized continuous problem (well posed under
\crefrange{H:mesh}{H:spaces}, and assumed smooth enough by \cref{H:regularity},
see \cref{sec:convergence}) one has $\mathcal{L}_{\ba}U=F$ in $L^2(K)$ for every $K$, and
therefore the method is consistent:
\begin{equation}\label{eq:consistency}
  \BASGS (\ba; U-U_h, V_h) = 0 \qquad \forall\, V_h\in\Xh .
\end{equation}

The stabilization parameters $\tau_{1,K}$, $\tau_{2,K}$ and $\tau_{1,\mathrm{NS}}$ are those of
\cref{eq:Tau1Final,eq:Tau2Final,eq:TauNavierStokes}, used here without restatement, with the
elemental evaluations fixed there: $c_1,c_2>0$ are the algorithmic constants,
$|\ba|_{\infty,K}\coloneqq\sup_K|\ba|$ (Euclidean norm) is the value at which the advection
modulus is taken on $K$, and $\alpha_K\coloneqq\alpha_{\infty,K}\coloneqq\sup_K\alpha$; we
also write $\alpha_{0,K}\coloneqq\inf_K\alpha>0$. What organizes the estimates below is not
$\tau_{1,\mathrm{NS}}$ itself but its porosity{\hyp}weighted inverse, which we abbreviate---in
analogy with~\cite{codina2001stabilized}---as
\begin{equation}\label{eq:phi1}
  \varphi_{1,K} \coloneqq \alpha_K\,\tau_{1,\mathrm{NS}}^{-1}
  = \alpha_K\Bigl(c_1\frac{\nu}{h_K^2}+c_2\frac{|\ba|_{\infty,K}}{h_K}\Bigr),
\end{equation}
so that $\tau_{1,K} = \bigl(\varphi_{1,K}+\sigma\bigr)^{-1}$. The damped reaction coefficient
$\sigmatilde$ of \cref{eq:SigmaAlpha} then admits, on each element, the equivalent forms
\begin{equation}\label{eq:sigmatilde}
  \sigmatilde\big|_K
  \coloneqq \frac{\tau_{1,\mathrm{NS}}^{-1}\sigma}{\tau_{1,\mathrm{NS}}^{-1}
  + \sigma/\alpha_K}
  = \frac{\varphi_{1,K}\,\sigma}{\varphi_{1,K}+\sigma}
  = \sigma - \sigma^2\tau_{1,K}
  = \sigma\,\varphi_{1,K}\,\tau_{1,K}\;\ge 0 ,
\end{equation}
of which the last two carry the whole reactive bookkeeping below. We drop the element subscript
($\tau_1$, $\tau_2$, $\varphi_1$, $\alpha_K$, $h$, $|\ba|_\infty$) whenever the elementwise
character is clear from the context.

\subsection{Standing assumptions and notation}\label{sec:assumptions}
Throughout, $C$ denotes a generic positive constant, possibly different at each occurrence,
that may depend on $d$, on the polynomial degrees, on the shape regularity of the mesh
family, on the algorithmic constants $c_1,c_2$ and on the constants
$C_{\nabla\alpha},C_2,c_J,c_J'$ named below, but \emph{not} on the mesh size,
nor on $\nu,\sigma,\varepsilon,\ba$ or $\alpha$.

The analysis is carried out under the standing assumptions of \cref{sec:StabilityASGS}: the
assumptions \cref{H:mesh,H:data,H:projector,H:porosity,H:advection,H:spaces,H:regularity} and,
specific to the ASGS method, the coercivity condition \cref{H:coercivity}
($c_1>2\xi\Cinva^2$ for some $\xi>2$) and---only when $\sigma>0$---the interior{\hyp}face jump
condition \cref{H:jump}. Each assumption is recalled
where it is first used, and \cref{rem:hypotheses} collects the bookkeeping for the continuity
proof. Three consequences are used so often that we record them here: \cref{H:porosity}
supplies $\delta_\alpha\coloneqq 1+C_{\nabla\alpha}$ with
$\alpha_{0,K}\le\alpha_K\le\delta_\alpha\,\alpha_{0,K}$ and $\alpha_K\le1$; \cref{H:advection}
gives $\nabla\cdot(\alpha\ba)=0$ in the weak sense, with $\alpha\ba$ continuous across interior
faces, hence a.e.\ in each $K$ (the form used in \cref{eq:elemibp}); and \cref{H:spaces} makes
$\tau_1,\tau_2$---and therefore $\sigmatilde$ through \cref{eq:sigmatilde}---constant on each
element, which is what allows them to be moved in and out of elemental integrals.

\begin{remark}[The porosity{\hyp}resolution condition and the constant $\delta_\alpha$]
\label{rem:porosityresolution}
Three readings of \cref{H:porosity} are used repeatedly below, and are worth separating
once. (i)~The porosity entering the variational forms is the exact field $\alpha$; it is not
interpolated onto a finite element space (that variant is left to future work, cf.\
\cref{sec:Conclusions}), so \cref{eq:resolved} is a condition on $\alpha$ and on the mesh alone.
Only the stabilization parameters replace $\alpha$ by the elementwise representative $\alpha_K$.
(ii)~The comparability $\alpha_{\infty,K}\le\delta_\alpha\,\alpha_{0,K}$, with
$\delta_\alpha=1+C_{\nabla\alpha}$, is a consequence of \cref{eq:resolved} and not a further
hypothesis, provided the elements are \emph{convex} and the length $h_K$ in \cref{eq:resolved}
is the diameter of $K$: any two points of $K$ are then joined by a segment lying in $K$, of
length at most $h_K$, whence
$\alpha_{\infty,K}-\alpha_{0,K}\le h_K\norm{\nabla\alpha}_{L^\infty(K)}\le
C_{\nabla\alpha}\,\alpha_{0,K}$. Convexity holds for every element family used
here---simplices and non{\hyp}degenerate parallelotopes---and is assumed throughout; on a
non{\hyp}convex element the factor becomes $1+C_{\mathrm{geo}}C_{\nabla\alpha}$, with
$C_{\mathrm{geo}}$ a uniform bound on the ratio of the interior path{\hyp}length metric of $K$
to its diameter. (iii)~If the length in \cref{eq:resolved} is not the diameter but a smaller
element{\hyp}size surrogate $\chi_1\le h_K$ with $h_K\le\chi_0\,\chi_1$ for a fixed
$\chi_0$---the convention in force when the stabilization parameters are built from a continuous
element{\hyp}size field rather than the raw diameter---the same argument yields the factor
$1+\chi_0\,C_{\nabla\alpha}$ instead.
\end{remark}

\subsection{Technical results}\label{sec:auxiliary}

We repeatedly use the standard inverse estimate: on a non{\hyp}degenerate mesh family there is a
constant $\Cinv$, independent of the mesh size, such that for every $\psi_h$ in the finite
element space restricted to $K$---that is, in the image on $K$ of a fixed finite{\hyp}dimensional
space on the reference element---
\begin{equation}
    \|\psi_h \|_{W^l_p(K)} \le \Cinv\, h^{m-l + d/p - d/q}\, \| \psi_h \|_{W^m_q(K)},
    \label{eq:InverseEstimateFiniteOrderNorm}
\end{equation}
valid for $0 \le m \le l$ and $1\le p,q \le \infty$ (see,
e.g.,~\cite[Thm.~4.5.11]{brenner2008mathematical}). We use it in its gradient form
($l=1$, $m=0$, $p=q=2$: $\lVert\nabla\psi_h\rVert_K\le\Cinv h^{-1}\norm{\psi_h}_K$) throughout,
and in its $L^\infty$ form ($l=m=0$, $p=\infty$, $q=2$:
$\norm{\psi_h}_{L^\infty(K)}\le\Cinv h^{-d/2}\norm{\psi_h}_K$) only in the jump terms of the
continuity proof.

\subsubsection{Parameter inequalities}

We first collect the elementary relations among $\tau_1,\tau_2,\varphi_1$ and $\sigmatilde$ that
are used, silently, throughout the coercivity and continuity proofs below.

\begin{lemma}[Parameter inequalities]\label{lem:parameters}
On each element $K$ (subscripts omitted) there hold:
\begin{enumerate}[label=\textup{(P\arabic*)},leftmargin=3.2em]
  \item\label{P:basic}
  $\sigma\tau_1\le 1$, $\varphi_1\tau_1\le 1$, $\tau_1\le \varphi_1^{-1}$;
  \item\label{P:components}
  $c_1\,\nu\,\alpha_K\le \varphi_1 h^2$ and
  $c_2\,\alpha_K\,|\ba|_{\infty}\,h\le \varphi_1 h^2$; in particular
  $\nu\,\tau_1\,\alpha_K\le h^2/c_1$;
  \item\label{P:27}
  $\varphi_1 h^2 = c_1\,\alpha_K^2\,\tau_2$; in particular
  $\varphi_1^{1/2} h = c_1^{1/2}\,\alpha_K\,\tau_2^{1/2}
  \le c_1^{1/2}\,\tau_2^{1/2}$;
  \item\label{P:sigmatilde}
  $\sigmatilde\le\min\{\sigma,\varphi_1\}$ and
  $\sigmatilde\,\tau_1\le 1$;
  \item\label{P:eps}
  under \eqref{eq:UpperBoundOnEpsilon}, $\varepsilon\tau_2\le C_2\,\varphi_1\tau_1\le C_2<1$;
  moreover, $\varepsilon h^2\le C_2\,c_1\,\alpha_K^2\,\tau_1$, and hence
  $\varepsilon^{1/2}\le c_1^{1/2}\,\alpha_K\,\tau_1^{1/2}/h$.
\end{enumerate}
\end{lemma}

\begin{proof}
Items \ref{P:basic} and \ref{P:components} are immediate from
\eqref{eq:phi1} and $\tau_1^{-1}=\varphi_1+\sigma$.

For \ref{P:27}, \eqref{eq:phi1} and \cref{eq:Tau2Final} give
\begin{equation*}
  \varphi_1 h^2 = \alpha_K\,h^2\,\tau_{1,\mathrm{NS}}^{-1}
  = c_1\,\alpha_K^2\,\frac{h^2}{c_1\,\alpha_K\,\tau_{1,\mathrm{NS}}}
  = c_1\,\alpha_K^2\,\tau_2 ,
\end{equation*}
and $\alpha_K\le1$ \cref{H:porosity} gives the second statement.

For \ref{P:sigmatilde}, \eqref{eq:sigmatilde} exhibits
$\sigmatilde=\sigma\varphi_1/(\varphi_1+\sigma)$ as a harmonic-type mean, hence bounded by
each of its two factors; moreover
\begin{equation*}
  \sigmatilde\,\tau_1 = \sigma\,\varphi_1\,\tau_1^2
  \le (\sigma\tau_1)\,(\varphi_1\tau_1)\le 1 .
\end{equation*}

For \ref{P:eps}, \eqref{eq:UpperBoundOnEpsilon} on the element $K$ reads
$\varepsilon\le C_2\,c_1\,\alpha_K^2\,\tau_1/h^2$; multiplying by $\tau_2$ and using
\ref{P:27},
\begin{equation*}
  \varepsilon\,\tau_2 \le C_2\,c_1\,\alpha_K^2\,\tau_1\,\frac{\tau_2}{h^2}
  = C_2\,\varphi_1\,\tau_1 \le C_2 .
\end{equation*}
The remaining statements are \eqref{eq:UpperBoundOnEpsilon} itself, followed by taking square
roots and using $C_2<1$.
\end{proof}

\subsubsection{Weighted inverse estimates}

Next, the porosity{\hyp}weighted forms of \eqref{eq:InverseEstimateFiniteOrderNorm}. They are
what the porous setting adds to the classical argument: $\alpha$ is not piecewise polynomial, so
the inverse estimate cannot be applied to the products $\alpha\ViscProj\nabla\boldsymbol{w}_h$
and $\alpha\boldsymbol{w}_h$ directly.

\begin{lemma}[Weighted inverse estimates]\label{lem:winv}
Let \cref{H:mesh,H:projector,H:porosity} hold. There is a constant
$\Cinva = \Cinva(\Cinv,C_{\nabla\alpha},d)$ such that, for all
finite element functions $\boldsymbol{w}_h\in\mathcal{V}_{h0}$,
$r_h\in\mathcal{Q}_{h0}$, and all $K\in\Th$,
\begin{align}
  \bigl\lVert\nabla\cdot\bigl(\alpha\ViscProj\nabla\boldsymbol{w}_h\bigr)
  \bigr\rVert_{K}
  &\le \frac{\Cinva}{h_K}\,\alpha_K^{1/2}\,
  \bigl\lVert\alpha^{1/2}\ViscProj\nabla\boldsymbol{w}_h\bigr\rVert_{K},
  \label{eq:winv-divvisc}\\
  \bigl\lVert\alpha^{1/2}\ViscProj\nabla\boldsymbol{w}_h\bigr\rVert_{K}
  &\le \frac{\Cinv}{h_K}\,\alpha_K^{1/2}\,
  \norm{\boldsymbol{w}_h}_{K},
  \label{eq:winv-grad}\\
  \bigl\lVert\nabla\cdot(\alpha\boldsymbol{w}_h)\bigr\rVert_{K}
  &\le \frac{\Cinva}{h_K}\,\alpha_K\,\norm{\boldsymbol{w}_h}_{K},
  \label{eq:winv-divu}\\
  \norm{\alpha\,\ba\cdot\nabla\boldsymbol{w}_h}_{K}
  &\le \frac{\Cinv}{h_K}\,\alpha_K\,|\ba|_{\infty,K}\,
  \norm{\boldsymbol{w}_h}_{K},
  \label{eq:winv-conv}\\
  \norm{\alpha\nabla r_h}_{K}
  &\le \frac{\Cinv}{h_K}\,\alpha_K\,\norm{r_h}_{K}.
  \label{eq:winv-gradp}
\end{align}
\end{lemma}

\begin{proof}
Write $\mathbf{M}\coloneqq\ViscProj\nabla\boldsymbol{w}_h$; since $\ViscProj$ acts pointwise with
constant coefficients, $\mathbf{M}$ inherits from $\boldsymbol{w}_h$ the finite element structure
on $K$ and \eqref{eq:InverseEstimateFiniteOrderNorm} applies to it. By the product
rule, $\nabla\cdot(\alpha\mathbf{M})=\alpha\,\nabla\cdot\mathbf{M}
+\mathbf{M}\,\nabla\alpha$. For the first contribution, using
$\norm{\nabla\cdot\mathbf{M}}_K\le\sqrt{d}\,\norm{\nabla\mathbf{M}}_K$, the inverse
estimate \eqref{eq:InverseEstimateFiniteOrderNorm}, and
$\norm{\mathbf{M}}_K\le\alpha_{0,K}^{-1/2}\norm{\alpha^{1/2}\mathbf{M}}_K$,
\[
\begin{aligned}
  \norm{\alpha\,\nabla\cdot\mathbf{M}}_K
  &\le \alpha_{\infty,K}\,\sqrt{d}\,\frac{\Cinv}{h_K}\,\norm{\mathbf{M}}_K \\
  &\le \sqrt{d}\,\frac{\Cinv}{h_K}\,\frac{\alpha_{\infty,K}}{\alpha_{0,K}^{1/2}}\,\norm{\alpha^{1/2}\mathbf{M}}_K \\
  &\le \sqrt{d\,\delta_\alpha}\,\frac{\Cinv}{h_K}\,\alpha_K^{1/2}\,\norm{\alpha^{1/2}\mathbf{M}}_K,
\end{aligned}
\]
since $\alpha_{\infty,K}/\alpha_{0,K}^{1/2}
=(\alpha_{\infty,K}/\alpha_{0,K})^{1/2}\alpha_{\infty,K}^{1/2}
\le\delta_\alpha^{1/2}\alpha_K^{1/2}$ by \cref{H:porosity}. For the second contribution,
by \eqref{eq:resolved},
\[
\begin{aligned}
  \norm{\mathbf{M}\,\nabla\alpha}_K
  &\le \norm{\nabla\alpha}_{L^\infty(K)}\norm{\mathbf{M}}_K \\
  &\le \frac{C_{\nabla\alpha}\,\alpha_{0,K}}{h_K}\,
  \alpha_{0,K}^{-1/2}\norm{\alpha^{1/2}\mathbf{M}}_K \\
  &\le \frac{C_{\nabla\alpha}}{h_K}\,\alpha_K^{1/2}\,
  \norm{\alpha^{1/2}\mathbf{M}}_K .
\end{aligned}
\]
This proves \eqref{eq:winv-divvisc} with
$\Cinva\coloneqq\sqrt{d\,\delta_\alpha}\,\Cinv+C_{\nabla\alpha}$. Estimate
\eqref{eq:winv-grad} follows from $\alpha\le\alpha_K$ on $K$, the pointwise contraction
property $|\ViscProj\nabla\boldsymbol{w}_h|\le|\nabla\boldsymbol{w}_h|$ and
\eqref{eq:InverseEstimateFiniteOrderNorm}. For \eqref{eq:winv-divu}, expand
$\nabla\cdot(\alpha\boldsymbol{w}_h)=\alpha\nabla\cdot\boldsymbol{w}_h
+\nabla\alpha\cdot\boldsymbol{w}_h$ and argue as above (using $\alpha_{0,K}\le\alpha_K$).
The estimates \eqref{eq:winv-conv} and \eqref{eq:winv-gradp} are immediate from
$\alpha\le\alpha_K$, $|\ba|\le|\ba|_{\infty,K}$ and \eqref{eq:InverseEstimateFiniteOrderNorm}.
\end{proof}

\begin{remark}\label{rem:winvconst}
\Cref{lem:winv} makes precise the inverse estimates used in the stability proof for the
$\alpha$-weighted second-order term $\nabla\cdot(\alpha\nu\ViscProj\nabla\buh)$: since
$\alpha$ is \emph{not} piecewise polynomial, \cref{eq:InverseEstimateFiniteOrderNorm} does not
apply to it directly, and the relevant constant is not the bare $\Cinv$ but the slightly larger
$\Cinva$, which incorporates the porosity-resolution constant $C_{\nabla\alpha}$ of
\eqref{eq:resolved}. The main text writes $\Cinv$ for this enlarged constant, by an
explicit abuse of notation; its conditions on $c_1$---in particular
\cref{eq:conditions_on_num_param} and the coercivity condition \cref{H:coercivity} of
\cref{lemma:Stability}---are therefore to be read with the enlarged constant already in place,
and require no further substitution. We write $\Cinva$ explicitly throughout this appendix
only to keep the two constants distinguishable within the proofs.
\end{remark}

\subsubsection{The jump of $\tau_1$}

The last elemental ingredient is needed only by the continuity proof, and there only when
$\sigma>0$: it controls how much $\tau_1$ may vary across an interior face. Note that, by the
definition \eqref{eq:phi1} of $\varphi_1$, the jump condition \cref{eq:JumpCondition} of
\cref{H:jump} reads simply
$c_J\,\varphi_{1,K_1}\le\varphi_{1,K_2}\le c_J'\,\varphi_{1,K_1}$ for any two elements sharing a
face, and it is in this form that we use it.

\begin{lemma}[Jump of $\tau_1$]\label{lem:jump}
Let \cref{H:jump} hold and let $K_1,K_2$ share the face $\Gamma^b$. Setting
$\jump{\tau_1}\coloneqq\tau_{1,K_1}-\tau_{1,K_2}$, there holds, for $i\in\{1,2\}$,
\begin{equation}\label{eq:jumpest}
  \sigma\,\bigl|\jump{\tau_1}\bigr|
  \le C\,\frac{\sigma\,\varphi_{1,K_i}}{(\varphi_{1,K_i}+\sigma)^2}
  = C\,\sigmatilde\big|_{K_i}\,\tau_{1,K_i}.
\end{equation}
\end{lemma}

\begin{proof}
We have
\begin{equation*}
  \jump{\tau_1}
  = \frac{\varphi_{1,K_2}-\varphi_{1,K_1}}
  {(\varphi_{1,K_1}+\sigma)(\varphi_{1,K_2}+\sigma)} ,
\end{equation*}
and \eqref{eq:JumpCondition} gives
$|\varphi_{1,K_1}-\varphi_{1,K_2}|\le
\max\{c_J'-1,\,1-c_J\}\,\varphi_{1,K_i}$ as well as
$\varphi_{1,K_2}+\sigma\ge \min\{c_J,1\}\,(\varphi_{1,K_1}+\sigma)$ (and symmetrically),
from which \eqref{eq:jumpest} follows with the identity \eqref{eq:sigmatilde}.
\end{proof}

\subsubsection{Two global integration-by-parts identities}

Both proofs below open by transferring a derivative from one argument to the other. The two
transfers that involve no boundary term---because the velocities vanish on $\Gamma$---are
recorded once here.

\begin{lemma}[Global identities]\label{lem:globalid}
Let \cref{H:porosity,H:advection} hold. For all $\bu,\bv\in H^1_0(\Omega)^d$ and all
$p,q\in H^1(\Omega)$,
\begin{align}
  (\bv,\alpha\,\ba\cdot\nabla\bu) &= -(\bu,\alpha\,\ba\cdot\nabla\bv), \label{eq:skew}\\
  (\bv,\alpha\nabla p) &= -\bigl(p,\nabla\cdot(\alpha\bv)\bigr),
  \qquad
  \bigl(q,\nabla\cdot(\alpha\bu)\bigr) = -(\bu,\alpha\nabla q). \label{eq:globalibp}
\end{align}
\end{lemma}

\begin{proof}
For \eqref{eq:skew}, the pointwise identity
$(\alpha\ba\cdot\nabla\bu)\cdot\bv+(\alpha\ba\cdot\nabla\bv)\cdot\bu
=\alpha\ba\cdot\nabla(\bu\cdot\bv)$ holds and $\bu\cdot\bv\in W^{1,1}_0(\Omega)$; since
$\alpha\ba\in L^\infty(\Omega)^d$ and $\nabla\cdot(\alpha\ba)=0$ in the weak sense
(\cref{H:advection}), $\int_\Omega\alpha\ba\cdot\nabla(\bu\cdot\bv)\dd\Omega=0$, which is
\eqref{eq:skew}. For \eqref{eq:globalibp}, $\alpha\in W^{1,\infty}(\Omega)$
(\cref{H:porosity}) and $\bv\in H^1_0(\Omega)^d$ give $\alpha\bv\in H^1_0(\Omega)^d$, and
likewise for $\bu$; the two identities are then the divergence theorem applied to
$p\,\alpha\bv$ and to $q\,\alpha\bu$, whose boundary terms vanish.
\end{proof}

\noindent Taking $\bv=\bu$ and $q=p$ in \cref{lem:globalid} gives the two \emph{diagonal}
identities used in the stability proof: the convective pairing $(\bu,\alpha\,\ba\cdot\nabla\bu)$
vanishes, and $(\bu,\alpha\nabla p)+\bigl(p,\nabla\cdot(\alpha\bu)\bigr)=0$.

\subsubsection{The ASGS triple norm}\label{sec:norm}

The stability and convergence analysis is carried out in the mesh{\hyp}dependent ASGS working
norm $\triplenorm{\cdot}_{\mathrm{A}}$ of \cref{lemma:Stability}; as this appendix treats only
the ASGS method, no ambiguity arises and we drop the subscript, writing simply
$\triplenorm{\cdot}$. Written out in the expanded broken notation of
\cref{eq:TauInnerProduct}, it reads
\begin{equation}\label{eq:triplenorm}
\begin{aligned}
  \triplenorm{V_h}
  \coloneqq \Bigl(
  &\nu\bigl\lVert\alpha^{1/2}\ViscProj\nabla\bvh\bigr\rVert^2
  + \bigl\lVert\sigmatilde^{1/2}\bvh\bigr\rVert_h^2
  + \varepsilon\bigl\lVert q_h \bigr\rVert^2 \\
  &+ \bigl\lVert\tau_1^{1/2}\alpha\X(V_h)\bigr\rVert_h^2
  + \bigl\lVert\tau_2^{1/2}\nabla\cdot(\alpha\bvh)\bigr\rVert_h^2
  \Bigr)^{1/2}.
\end{aligned}
\end{equation}
Its five terms control, respectively, the viscous gradient, the damped reaction, the
compressibility, the first{\hyp}order momentum residual and the porous divergence. That it is
indeed a norm on $\Xh$---the viscous term alone controlling the velocity, after which the
$\tau_1$ residual term and the compressibility (or, when $\varepsilon=0$, the zero{\hyp}mean
constraint) pin the pressure---is the content of the following lemma.

\begin{lemma}[Definiteness of the working norm]\label{lem:definiteness}
Let \cref{H:data,H:projector,H:porosity,H:advection,H:spaces} hold (recall that
\cref{H:spaces} fixes the pressure constant by the zero{\hyp}mean normalization when
$\varepsilon = 0$). Then $\triplenorm{\cdot}$ of
\eqref{eq:triplenorm} is a norm on $\Xh$.
\end{lemma}
\begin{proof}
Homogeneity, symmetry and the triangle inequality being immediate, only definiteness remains.
Suppose $\triplenorm{V_h}=0$. Since $\nu>0$ (\cref{H:data}) and $\alpha\ge\alpha_0>0$
(\cref{H:porosity}), the first term forces $\ViscProj\nabla\bvh=0$, hence $\nabla\bvh=0$ by the
Korn compatibility \cref{eq:PiKorn} of \cref{H:projector}---available with
$C_{\mathrm{K}}=\sqrt2$ for the whole family of \cref{lem:projinstances}, by
\cref{lem:projfamily}---and $\bvh=\boldsymbol 0$ by Poincar\'e. Then $\X(V_h)=\nabla p_h$, and
the $\tau_1$ term---whose weights $\tau_{1,K}=(\varphi_{1,K}+\sigma)^{-1}$ are positive, $|\ba|_{\infty,K}$ being finite by \cref{H:advection}---forces $\nabla p_h=0$, so that
$p_h$ is constant on the connected domain $\Omega$ ($\mathcal{Q}_{h0}\subset H^1(\Omega)$ by
\cref{H:spaces}). That constant vanishes: through the term $\varepsilon\norm{p_h}^2$ when
$\varepsilon>0$, and through the zero{\hyp}mean constraint defining $\mathcal{Q}_0$ when
$\varepsilon=0$.
\end{proof}

\subsection{Stability}\label{sec:coercivity}
The stability of the ASGS method is the coercivity of $\BASGS $ in the working norm
\eqref{eq:triplenorm}: this is \cref{lemma:Stability} of the main text, which we now prove. It
is the only place where the design condition \cref{H:coercivity} on $c_1$ is used, and the first
of the two where the porosity{\hyp}weighted inverse estimate \eqref{eq:winv-divvisc} is (the
other being Step~4 of the continuity proof).

\begin{lemma}[Coercivity]\label{lem:coercivity}
Let \cref{H:mesh,H:data,H:projector,H:porosity,H:advection,H:spaces} hold, together with the design
condition \cref{H:coercivity}, i.e.\ $c_1 > 2\xi\,\Cinva^2$ for some $\xi>2$. Then
\begin{equation}\label{eq:coercivity}
  \BASGS (\ba; U_h,U_h) \ \ge\ C_{\mathrm{coer}}\,\triplenorm{U_h}^2
  \qquad\forall\,U_h\in\Xh ,
\end{equation}
with the explicit constant
\begin{equation}\label{eq:coerconstant}
  C_{\mathrm{coer}} \coloneqq \min\Bigl\{1-\frac{4\Cinva^2}{c_1},\ 1-C_2\Bigr\}
  \ \ge\ \min\Bigl\{1-\frac{2}{\xi},\ 1-C_2\Bigr\} \ >\ 0 ,
\end{equation}
$C_2<1$ being the compressibility constant of \eqref{eq:UpperBoundOnEpsilon}. In particular
$C_{\mathrm{coer}}$ is bounded below by a quantity depending only on $\xi$ and $C_2$, and is
independent of the mesh size and of $\nu,\sigma,\varepsilon,\ba,\alpha$.
\end{lemma}

\begin{proof}
Put $V_h=U_h=[\buh;p_h]$ in \eqref{eq:Bstab} and group the terms.

\medskip\noindent\emph{Step 1: the Galerkin diagonal.}
By the diagonal case of \cref{lem:globalid}, the convective pairing vanishes,
$\bigl(\buh,\alpha\,\ba\cdot\nabla\buh\bigr)=0$ by \eqref{eq:skew}, and the two
pressure--velocity pairings cancel,
$\bigl(\buh,\alpha\nabla p_h\bigr)+\bigl(p_h,\nabla\cdot(\alpha\buh)\bigr)=0$ by
\eqref{eq:globalibp}. Hence the Galerkin part reduces to the coercive diagonal of the main
text,
$2\nu\norm{\alpha^{1/2}\ViscProj\nabla\buh}^2+\sigma\norm{\buh}^2+\varepsilon\norm{p_h}^2$.

\medskip\noindent\emph{Step 2: the stabilization diagonals.}
In the $\tau_1$ term the mixed products cancel by the symmetry of the weighted inner
product---the weights $\tau_{1,K}$ being elementwise constant, \cref{H:spaces}---leaving
\begin{equation}\label{eq:stabdiag1}
  \bigl(\alpha\X(U_h)+G(\buh),\,\tau_1[\alpha\X(U_h)-G(\buh)]\bigr)_h
  = \bigl\lVert\tau_1^{1/2}\alpha\X(U_h)\bigr\rVert_h^2
  - \bigl\lVert\tau_1^{1/2}G(\buh)\bigr\rVert_h^2 ,
\end{equation}
and, likewise, the $\tau_2$ term equals
$\lVert\tau_2^{1/2}\nabla\cdot(\alpha\buh)\rVert_h^2
-\varepsilon^2\lVert\tau_2^{1/2}p_h\rVert_h^2$.

\medskip\noindent\emph{Step 3: reaction renormalization and penalty absorption.}
Write $G(\buh)=G_{\mathrm v}-\sigma\buh$ with
$G_{\mathrm v}\coloneqq 2\nu\,\nabla\cdot(\alpha\ViscProj\nabla\buh)$. Then
\begin{equation}\label{eq:Gexpand}
  \bigl\lVert\tau_1^{1/2}G(\buh)\bigr\rVert_h^2
  = \bigl\lVert\tau_1^{1/2}G_{\mathrm v}\bigr\rVert_h^2
  - 2\sigma\bigl(\tau_1 G_{\mathrm v},\buh\bigr)_h
  + \sigma^2\bigl\lVert\tau_1^{1/2}\buh\bigr\rVert_h^2 ,
\end{equation}
and its last term recombines with the Galerkin reaction into the damped coefficient of the
working norm: on each element $\sigma-\sigma^2\tau_{1,K}=\sigmatilde|_K$ by
\eqref{eq:sigmatilde}, so that, summing (again using \cref{H:spaces}),
\begin{equation}\label{eq:reactrenorm}
  \sigma\norm{\buh}^2-\sigma^2\bigl\lVert\tau_1^{1/2}\buh\bigr\rVert_h^2
  = \bigl\lVert\sigmatilde^{1/2}\buh\bigr\rVert_h^2 .
\end{equation}
Similarly, by \ref{P:eps} ($\varepsilon\tau_2\le C_2<1$) the penalty residual is absorbed,
$\varepsilon\norm{p_h}^2-\varepsilon^2\lVert\tau_2^{1/2}p_h\rVert_h^2
\ge (1-C_2)\,\varepsilon\norm{p_h}^2\ge 0$.

\medskip\noindent\emph{Step 4: absorption of the viscous residual.}
By \eqref{eq:winv-divvisc} and $\tau_1\le h^2/(c_1\alpha_K\nu)$ from \ref{P:components},
\begin{equation}\label{eq:viscabsorb}
  \bigl\lVert\tau_1^{1/2}G_{\mathrm v}\bigr\rVert_h^2
  = 4\nu^2\sum_K\tau_{1,K}\bigl\lVert\nabla\cdot(\alpha\ViscProj\nabla\buh)\bigr\rVert_K^2
  \le \frac{4\,\Cinva^2}{c_1}\;\nu\,\norm{\alpha^{1/2}\ViscProj\nabla\buh}^2 ,
\end{equation}
and the cross term is estimated per element by the Cauchy--Schwarz inequality, the inverse
estimate \eqref{eq:winv-divvisc}, and Young's inequality $2ab\le\eta^{-1}a^2+\eta b^2$ (any
$\eta>0$): since $G_{\mathrm v}=2\nu\,\nabla\cdot(\alpha\ViscProj\nabla\buh)$,
\[
\begin{aligned}
  \bigl|2\sigma\bigl(\tau_1 G_{\mathrm v},\buh\bigr)_K\bigr|
  &\le 4\nu\sigma\tau_1
    \bigl\lVert\nabla\cdot(\alpha\ViscProj\nabla\buh)\bigr\rVert_K\norm{\buh}_K \\
  &\le 4\nu\sigma\tau_1\,\frac{\Cinva}{h}\,\alpha_K^{1/2}
    \bigl\lVert\alpha^{1/2}\ViscProj\nabla\buh\bigr\rVert_K\norm{\buh}_K ,
\end{aligned}
\]
so that, taking $a=2(\nu\sigma\tau_1)^{1/2}\bigl\lVert\alpha^{1/2}\ViscProj\nabla\buh\bigr\rVert_K$
and $b=(\nu\sigma\tau_1)^{1/2}(\Cinva/h)\,\alpha_K^{1/2}\norm{\buh}_K$,
\begin{equation}\label{eq:crossbound}
  2\sigma\bigl(\tau_1 G_{\mathrm v},\buh\bigr)_h
  \ \ge\ -\sum_K\Bigl\{\tfrac{4}{\eta}\,\nu\sigma\tau_1\bigl\lVert\alpha^{1/2}\ViscProj\nabla\buh\bigr\rVert_K^2
  + \eta\,\Cinva^2\,\frac{\nu\sigma\tau_1\alpha_K}{h^2}\,\norm{\buh}_K^2\Bigr\}.
\end{equation}

\medskip\noindent\emph{Step 5: collection.}
We combine the diagonal, the reaction renormalization \eqref{eq:reactrenorm}, the penalty
absorption, and \eqref{eq:viscabsorb}--\eqref{eq:crossbound}, bounding $\sigma\tau_1\le1$
(\ref{P:basic}) in the first cross term of \eqref{eq:crossbound} and, in the second,
$c_1\nu\alpha_K/h^2\le\varphi_1$ (\ref{P:components}) together with
$\sigmatilde=\sigma\varphi_1\tau_1$ to convert
$\nu\sigma\tau_1\alpha_K/h^2\le\sigmatilde/c_1$. Since both negative viscous contributions
subtract from the single Galerkin term $2\nu\norm{\alpha^{1/2}\ViscProj\nabla\buh}^2$, we obtain
\begin{equation}\label{eq:coerccollect}
\begin{aligned}
  \BASGS (\ba;U_h,U_h)
  \ \ge\ &\Bigl(2\bigl(1-\tfrac{2}{\eta}\bigr)-\tfrac{4\Cinva^2}{c_1}\Bigr)\,
  \nu\norm{\alpha^{1/2}\ViscProj\nabla\buh}^2
  + \Bigl(1-\tfrac{\eta\Cinva^2}{c_1}\Bigr)\bigl\lVert\sigmatilde^{1/2}\buh\bigr\rVert_h^2 \\
  &+ (1-C_2)\,\varepsilon\norm{p_h}^2
  + \bigl\lVert\tau_1^{1/2}\alpha\X(U_h)\bigr\rVert_h^2
  + \bigl\lVert\tau_2^{1/2}\nabla\cdot(\alpha\buh)\bigr\rVert_h^2 .
\end{aligned}
\end{equation}

\medskip\noindent\emph{Step 6: choice of the Young parameter, and conclusion.}
Only the first two coefficients of \eqref{eq:coerccollect} depend on $\eta>0$, which is still at
our disposal: writing $t\coloneqq\Cinva^2/c_1$, the viscous one,
$c_{\mathrm v}(\eta)\coloneqq2(1-2/\eta)-4t$, increases with $\eta$, while the damped{\hyp}reaction one,
$c_{\mathrm r}(\eta)\coloneqq1-\eta t$, decreases. Their minimum is therefore largest where the two
coincide, and $c_{\mathrm v}(\eta)=c_{\mathrm r}(\eta)$ reads $t\eta^2+(1-4t)\eta-4=0$, i.e.\
$(\eta-4)(t\eta+1)=0$: the optimal choice is $\eta=4$, \emph{whatever the value of $t$}, and it
makes both coefficients equal to
\begin{equation}\label{eq:coercoeff}
  c_{\mathrm v}(4)=c_{\mathrm r}(4)=1-\frac{4\Cinva^2}{c_1} .
\end{equation}
This common value is positive precisely when $c_1>4\Cinva^2$, so that condition is not merely
sufficient but necessary for the present collection---and it is exactly what
\cref{H:coercivity} grants, since ``$c_1>2\xi\Cinva^2$ for some $\xi>2$'' and
``$c_1>4\Cinva^2$'' are equivalent requirements on $c_1$. The first reading gives in addition
$4\Cinva^2/c_1<2/\xi$, whence the $\xi$-explicit lower bound in \eqref{eq:coerconstant}.

The compressibility coefficient is $1-C_2>0$ and the two residual squares carry the coefficient
$1$; comparing \eqref{eq:coerccollect} termwise with \eqref{eq:triplenorm} therefore yields
\eqref{eq:coercivity} with the constant \eqref{eq:coerconstant}.
\end{proof}

\noindent Because $\triplenorm{\cdot}$ is a \emph{norm} on $\Xh$ (\cref{lem:definiteness}) and
not merely a seminorm, \eqref{eq:coercivity} makes the operator of the square discrete system
injective, hence invertible: the discrete ASGS problem has a unique solution $U_h\in\Xh$, which
obeys the a priori bound
$\triplenorm{U_h}\le C_{\mathrm{coer}}^{-1}\sup_{V_h\in\Xh\setminus\{0\}}
\lvert\LASGS(\ba;V_h)\rvert/\triplenorm{V_h}$.

\begin{remark}[Sharpness of the condition on $c_1$]\label{rem:c1sharp}
The threshold $c_1>4\Cinva^2$ is the worst case over the reaction regime, and it can be traced
to the single place where the reaction was discarded: in Step~5 the factor $\sigma\tau_1\le1$
was used to charge the whole first cross term of \eqref{eq:crossbound} against the viscous
Galerkin term. Since Young's inequality is applied elementwise, $\eta$ may be taken
element{\hyp}dependent and $s_K\coloneqq\sigma\tau_{1,K}\in[0,1)$ kept: the elemental
coefficients are then $c_{\mathrm v,K}(\eta)=2-4s_K/\eta-4t$ and $c_{\mathrm r,K}(\eta)=1-\eta t$, whose maximin is
positive if and only if $2-4t(1+s_K)>0$, that is
\begin{equation}\label{eq:c1elemental}
  c_1 > 2\bigl(1+\sigma\tau_{1,K}\bigr)\,\Cinva^2 \qquad\text{on every } K .
\end{equation}
The reaction{\hyp}dominated limit $\sigma\tau_{1,K}\to1$ recovers $c_1>4\Cinva^2$, whereas in the
absence of reaction, and more generally on any viscosity{\hyp}dominated element,
$c_1>2\Cinva^2$ already suffices. The latter is the purely geometric, elementwise requirement
that governs the choice of $c_1$ in three dimensions
(cf.\ \cref{sec:NumericalExamplesStationaryCase3D}). We state \cref{lem:coercivity} under the
uniform condition \cref{H:coercivity} because it is parameter{\hyp}free.
\end{remark}

\subsection{Continuity}\label{sec:continuity}
The counterpart of \cref{lem:coercivity} is a continuity estimate in which the first argument of
$\BASGS $ is measured in the mesh{\hyp}dependent quantities that the interpolation error will
later supply. It is \cref{lemma:Continuity} of the main text. Of the design constants it
requires only $c_1,c_2>0$---not the coercivity condition \cref{H:coercivity}---but it does
require the jump condition \cref{H:jump}, and with it the $L^\infty$ inverse estimate, whenever
$\sigma>0$ (\cref{rem:hypotheses}).

\begin{lemma}[Continuity]\label{lem:continuity}
Let \cref{H:mesh,H:data,H:projector,H:porosity,H:advection,H:spaces} hold, together with the jump condition \cref{H:jump} when $\sigma>0$, with $c_1,c_2>0$. Then there exists $C>0$,
independent of the mesh size and of $\nu,\sigma,\varepsilon,\ba,\alpha$, such that
\begin{equation}\label{eq:continuity}
  \bigl\lvert\BASGS (\ba; U_h,V_h)\bigr\rvert
  \le C\Bigl(
  \bigl\lVert\tau_2^{1/2}h^{-1}\buh\bigr\rVert_h
  + \bigl\lVert\tau_1^{1/2}h^{-1}p_h\bigr\rVert_h
  \Bigr)\,\triplenorm{V_h}
  \qquad\forall\, U_h,V_h\in\Xh .
\end{equation}
Moreover, the proof yields the sharper bound
\begin{equation}\label{eq:sharpcont}
  \bigl\lvert\BASGS (\ba; U_h,V_h)\bigr\rvert
  \le C\Bigl(
  \bigl\lVert\alpha_K\,\tau_2^{1/2}h^{-1}\buh\bigr\rVert_h
  + \bigl\lVert\alpha_K\,\tau_1^{1/2}h^{-1}p_h\bigr\rVert_h
  \Bigr)\,\triplenorm{V_h},
\end{equation}
where $\alpha_K$ denotes the piecewise constant function equal to $\alpha_{\infty,K}$ on
each $K\in\Th$; \eqref{eq:continuity} follows from \eqref{eq:sharpcont} since
$\alpha\le 1$.
\end{lemma}

\begin{proof}
Throughout the proof we write $U=U_h=[\bu;p]$ and $V=V_h=[\bv;q]$, dropping the subscript
$h$ of the unknowns for conciseness, and, as elsewhere, we omit the element subscript inside
elementwise sums. The proof follows the steps of the proof of Lemma~2
in~\cite{codina2001stabilized}.

\medskip\noindent\emph{Step 0: term decomposition.}
Expanding the products in \eqref{eq:Bstab} we obtain
$\BASGS (\ba;U,V) = \sum_{j=1}^{18} T_j$, with
\begin{align}
  T_1 &\coloneqq 2\nu\bigl(\ViscProj\nabla\bv,\alpha\ViscProj\nabla\bu\bigr), \label{eq:T1}\\
  T_2 &\coloneqq \bigl(\bv,\alpha\X(U)\bigr), \label{eq:T2}\\
  T_3 &\coloneqq \sigma\,(\bv,\bu), \label{eq:T3}\\
  T_4 &\coloneqq \bigl(q,\nabla\cdot(\alpha\bu)\bigr), \label{eq:T4}\\
  T_5 &\coloneqq \varepsilon\,(q,p), \label{eq:T5}\\
  T_6 &\coloneqq -4\nu^2\bigl(\tau_1\nabla\cdot(\alpha\ViscProj\nabla\bv),
        \nabla\cdot(\alpha\ViscProj\nabla\bu)\bigr)_h, \label{eq:T6}\\
  T_7 &\coloneqq -2\nu\bigl(\alpha\X(V),
        \tau_1\nabla\cdot(\alpha\ViscProj\nabla\bu)\bigr)_h, \label{eq:T7}\\
  T_8 &\coloneqq 2\nu\sigma\bigl(\tau_1\bv,
        \nabla\cdot(\alpha\ViscProj\nabla\bu)\bigr)_h, \label{eq:T8}\\
  T_9 &\coloneqq 2\nu\bigl(\tau_1\nabla\cdot(\alpha\ViscProj\nabla\bv),
        \alpha\X(U)\bigr)_h, \label{eq:T9}\\
  T_{10} &\coloneqq \bigl(\alpha\X(V),\tau_1\alpha\X(U)\bigr)_h, \label{eq:T10}\\
  T_{11} &\coloneqq -\sigma\bigl(\tau_1\alpha\X(U),\bv\bigr)_h, \label{eq:T11}\\
  T_{12} &\coloneqq 2\nu\sigma\bigl(\tau_1\nabla\cdot(\alpha\ViscProj\nabla\bv),
        \bu\bigr)_h, \label{eq:T12}\\
  T_{13} &\coloneqq \sigma\bigl(\tau_1\bu,\alpha\X(V)\bigr)_h, \label{eq:T13}\\
  T_{14} &\coloneqq -\sigma^2\bigl(\tau_1\bu,\bv\bigr)_h, \label{eq:T14}\\
  T_{15} &\coloneqq -\varepsilon^2\bigl(\tau_2\,p,q\bigr)_h, \label{eq:T15}\\
  T_{16} &\coloneqq \varepsilon\bigl(\nabla\cdot(\alpha\bv),\tau_2\,p\bigr)_h, \label{eq:T16}\\
  T_{17} &\coloneqq -\varepsilon\bigl(\tau_2\nabla\cdot(\alpha\bu),q\bigr)_h, \label{eq:T17}\\
  T_{18} &\coloneqq \bigl(\nabla\cdot(\alpha\bv),\tau_2\nabla\cdot(\alpha\bu)\bigr)_h .
        \label{eq:T18}
\end{align}
Terms \eqref{eq:T1}--\eqref{eq:T5} come from the Galerkin part,
\eqref{eq:T6}--\eqref{eq:T14} from the momentum stabilization and
\eqref{eq:T15}--\eqref{eq:T18} from the mass stabilization in \eqref{eq:Bstab}; they
correspond, in this order, to terms (42)--(59) of~\cite{codina2001stabilized}.

\medskip\noindent\emph{Step 1: terms bounded directly by the working norm.}
By the Cauchy--Schwarz inequality (elementwise and then for the sums over elements),
\begin{align}
  |T_1| + |T_{10}| + |T_{18}|
  \le{}& 2\nu\bigl\lVert\alpha^{1/2}\ViscProj\nabla\bu\bigr\rVert
        \bigl\lVert\alpha^{1/2}\ViscProj\nabla\bv\bigr\rVert
  + \bigl\lVert\tau_1^{1/2}\alpha\X(U)\bigr\rVert_h
        \bigl\lVert\tau_1^{1/2}\alpha\X(V)\bigr\rVert_h \nonumber\\
  &+ \bigl\lVert\tau_2^{1/2}\nabla\cdot(\alpha\bu)\bigr\rVert_h
        \bigl\lVert\tau_2^{1/2}\nabla\cdot(\alpha\bv)\bigr\rVert_h
  \;\le\; 2\,\triplenorm{U}\,\triplenorm{V}.
  \label{eq:easystep}
\end{align}

\medskip\noindent\emph{Step 2: the reactive pair.}
Since $\tau_1$ is constant on each element and
$\sigma-\sigma^2\tau_{1}=\sigmatilde$ on each element by \eqref{eq:sigmatilde},
\begin{equation}\label{eq:reactivestep}
  T_3 + T_{14} = \bigl(\sigmatilde\,\bu,\bv\bigr)_h
  \le \bigl\lVert\sigmatilde^{1/2}\bu\bigr\rVert_h
      \bigl\lVert\sigmatilde^{1/2}\bv\bigr\rVert_h
  \le \triplenorm{U}\,\triplenorm{V}.
\end{equation}

\medskip\noindent\emph{Step 3: the penalty pair.}
By \cref{lem:parameters}\ref{P:eps}, $\varepsilon\tau_{2,K}\le C_2<1$ on every element,
hence
\begin{equation}\label{eq:penaltystep}
  |T_5 + T_{15}|
  \le \varepsilon\,\norm{p}\norm{q} + \varepsilon\,C_2\,\norm{p}\norm{q}
  \le 2\,\varepsilon^{1/2}\norm{p}\;\varepsilon^{1/2}\norm{q}
  \le 2\,\triplenorm{U}\,\triplenorm{V}.
\end{equation}

\medskip\noindent\emph{Step 4: terms containing the viscous operator.}
The key elemental estimate is obtained from \eqref{eq:winv-divvisc} together with
$\nu\tau_1\alpha_K\le h^2/c_1$ (\cref{lem:parameters}\ref{P:components}): for any finite
element velocity $\boldsymbol{w}_h$,
\begin{equation}\label{eq:keyvisc}
\begin{aligned}
  \bigl\lVert\tau_1^{1/2}\,2\nu\,
  \nabla\cdot(\alpha\ViscProj\nabla\boldsymbol{w}_h)\bigr\rVert_{K}
  &\le 2\nu\,\tau_1^{1/2}\,\frac{\Cinva}{h}\,\alpha_K^{1/2}\,
  \bigl\lVert\alpha^{1/2}\ViscProj\nabla\boldsymbol{w}_h\bigr\rVert_{K} \\
  &\le \frac{2\Cinva}{c_1^{1/2}}\;\nu^{1/2}
  \bigl\lVert\alpha^{1/2}\ViscProj\nabla\boldsymbol{w}_h\bigr\rVert_{K}.
\end{aligned}
\end{equation}
Consequently, by the Cauchy--Schwarz inequality and \eqref{eq:keyvisc} applied to $\bu$
and/or $\bv$,
\begin{align}
  |T_6| &\le
  \bigl\lVert\tau_1^{1/2}2\nu\nabla\cdot(\alpha\ViscProj\nabla\bu)\bigr\rVert_h
  \bigl\lVert\tau_1^{1/2}2\nu\nabla\cdot(\alpha\ViscProj\nabla\bv)\bigr\rVert_h \nonumber\\
  &\le \frac{4\Cinva^2}{c_1}\,
  \nu\bigl\lVert\alpha^{1/2}\ViscProj\nabla\bu\bigr\rVert
  \bigl\lVert\alpha^{1/2}\ViscProj\nabla\bv\bigr\rVert
  \le C\,\triplenorm{U}\triplenorm{V}, \label{eq:T6bound}\\
  |T_7| &\le
  \bigl\lVert\tau_1^{1/2}\alpha\X(V)\bigr\rVert_h
  \bigl\lVert\tau_1^{1/2}2\nu\nabla\cdot(\alpha\ViscProj\nabla\bu)\bigr\rVert_h \nonumber\\
  &\le \frac{2\Cinva}{c_1^{1/2}}\,
  \nu^{1/2}\bigl\lVert\alpha^{1/2}\ViscProj\nabla\bu\bigr\rVert\,
  \bigl\lVert\tau_1^{1/2}\alpha\X(V)\bigr\rVert_h
  \le C\,\triplenorm{U}\triplenorm{V}, \label{eq:T7bound}
\end{align}
and $|T_9|\le C\triplenorm{U}\triplenorm{V}$ in the same way as $T_7$, with the roles of
$\bu$ and $\bv$ interchanged. For $T_8$, the chained estimates
\eqref{eq:winv-divvisc}--\eqref{eq:winv-grad} give the ``double'' inverse estimate
\begin{equation}\label{eq:doubleinv}
  \bigl\lVert\nabla\cdot(\alpha\ViscProj\nabla\boldsymbol{w}_h)\bigr\rVert_{K}
  \le \frac{\Cinva\,\Cinv}{h^2}\,\alpha_K\,\norm{\boldsymbol{w}_h}_{K},
\end{equation}
whence, using $c_1\nu\alpha_K/h^2\le\varphi_1$ (\cref{lem:parameters}\ref{P:components})
and $\sigma\varphi_1\tau_1=\sigmatilde$,
\begin{align}
  |T_8|
  &\le \sum_{K} 2\nu\sigma\tau_1\,\frac{\Cinva\Cinv}{h^2}\,\alpha_K\,
  \norm{\bu}_K\norm{\bv}_K
  = \frac{2\Cinva\Cinv}{c_1}\sum_{K}
  \sigma\,\Bigl(\frac{c_1\nu\alpha_K}{h^2}\Bigr)\tau_1\,\norm{\bu}_K\norm{\bv}_K
  \nonumber\\
  &\le \frac{2\Cinva\Cinv}{c_1}
  \bigl\lVert\sigmatilde^{1/2}\bu\bigr\rVert_h
  \bigl\lVert\sigmatilde^{1/2}\bv\bigr\rVert_h
  \le C\,\triplenorm{U}\triplenorm{V},
  \label{eq:T8bound}
\end{align}
and $|T_{12}|\le C\triplenorm{U}\triplenorm{V}$ identically, applying
\eqref{eq:doubleinv} to $\bv$ instead of $\bu$. Summarizing this step,
\begin{equation}\label{eq:viscstep}
  |T_6|+|T_7|+|T_8|+|T_9|+|T_{12}| \le C\,\triplenorm{U}\,\triplenorm{V}.
\end{equation}

\medskip\noindent\emph{Step 5: splitting of $T_{13}$.}
Write $T_{13}=T_{13}^{\mathrm{c}}+R$, with the convective contribution
$T_{13}^{\mathrm{c}}\coloneqq\sigma(\tau_1\bu,\alpha\,\ba\cdot\nabla\bv)_h$ and
$R\coloneqq\sigma(\tau_1\bu,\alpha\nabla q)_h$. For $T_{13}^{\mathrm{c}}$,
\eqref{eq:winv-conv} and $c_2\alpha_K|\ba|_\infty/h\le\varphi_1$
(\cref{lem:parameters}\ref{P:components}) give
\begin{align}
  \bigl|\sigma(\tau_1\bu,\alpha\,\ba\cdot\nabla\bv)_h\bigr|
  &\le \sum_K \sigma\tau_1\,\alpha_K|\ba|_{\infty}\frac{\Cinv}{h}\,
  \norm{\bu}_K\norm{\bv}_K
  \le \frac{\Cinv}{c_2}\sum_K \sigma\varphi_1\tau_1\norm{\bu}_K\norm{\bv}_K
  \nonumber\\
  &\le \frac{\Cinv}{c_2}\,
  \bigl\lVert\sigmatilde^{1/2}\bu\bigr\rVert_h
  \bigl\lVert\sigmatilde^{1/2}\bv\bigr\rVert_h .
  \label{eq:T13conv}
\end{align}
The second contribution $R$ is kept and treated in the next step, where it is
absorbed into the group $N$.

\medskip\noindent\emph{Step 6: the convective--pressure group
$N\coloneqq R+T_2+T_4+T_{11}$.}

\smallskip\noindent\emph{Step 6a: exact rewriting of $N$.}
Besides the two global identities \eqref{eq:skew}--\eqref{eq:globalibp} of
\cref{lem:globalid}---applicable here because $\bu,\bv\in\mathcal{V}_{h0}\subset H^1_0(\Omega)^d$
and $p,q\in\mathcal{Q}_{h0}\subset H^1(\Omega)$---we need their elementwise counterpart for the
convective term, which does carry a boundary contribution: on each $K$, with $\tau_1$ constant
on $K$ (\cref{H:spaces}), $\ba$ continuous on $\overline{K}$ and $\nabla\cdot(\alpha\ba)=0$ in
$K$, the divergence theorem gives
\begin{equation}\label{eq:elemibp}
  \bigl(\tau_1\,\alpha\,\ba\cdot\nabla\bu,\bv\bigr)_K
  = -\bigl(\tau_1\,\alpha\,\ba\cdot\nabla\bv,\bu\bigr)_K
  + \tau_{1}\int_{\partial K}\alpha\,(\bn\cdot\ba)\,(\bu\cdot\bv)\dd\Gamma .
\end{equation}
Using \eqref{eq:skew}--\eqref{eq:globalibp} in $T_2$ and $T_4$,
\begin{equation*}
  T_2 + T_4 = -\bigl(\bu,\alpha\X(V)\bigr) - \bigl(p,\nabla\cdot(\alpha\bv)\bigr),
\end{equation*}
while \eqref{eq:elemibp} applied to the convective part of $T_{11}$ gives
\begin{equation*}
  T_{11}
  = \sigma\bigl(\tau_1\,\alpha\,\ba\cdot\nabla\bv,\bu\bigr)_h
  - \sigma\sum_{K}\tau_{1,K}\int_{\partial K}\alpha(\bn_K\cdot\ba)(\bu\cdot\bv)\dd\Gamma
  - \sigma\bigl(\tau_1\bv,\alpha\nabla p\bigr)_h .
\end{equation*}
Adding $R=\sigma(\tau_1\bu,\alpha\nabla q)_h$ and collecting terms,
\begin{align}
  N ={}& \underbrace{-\sigma\bigl(\tau_1\bv,\alpha\nabla p\bigr)_h}_{\mathrm{[I]}}
  \;+\; \underbrace{\sigma\bigl(\tau_1\bu,\alpha\X(V)\bigr)_h}_{\mathrm{[II]}}
  \;\underbrace{-\,\sigma\sum_{K}\tau_{1,K}\int_{\partial K}
  \alpha(\bn_K\cdot\ba)(\bu\cdot\bv)\dd\Gamma}_{\mathrm{[III]}}
  \nonumber\\
  &\;\underbrace{-\,\bigl(p,\nabla\cdot(\alpha\bv)\bigr)}_{\mathrm{[IV]}}
  \;\underbrace{-\,\bigl(\bu,\alpha\X(V)\bigr)}_{\mathrm{[V]}} ,
  \label{eq:fiveterms}
\end{align}
which is the porous counterpart of Eq.~(66) of~\cite{codina2001stabilized}. (Note that the
inter-element term [III] carries a \emph{minus} sign; in~\cite{codina2001stabilized} it is printed
with a plus sign, a harmless typo since the term is estimated in absolute value.)

\smallskip\noindent\emph{Step 6b: the pressure terms \textup{[I]+[IV]}.}
Integrating by parts elementwise, $(\tau_1\bv,\alpha\nabla p)_K
= -(\tau_1\nabla\cdot(\alpha\bv),p)_K
+ \tau_{1}\int_{\partial K}\alpha(\bn\cdot\bv)\,p\dd\Gamma$, and using
$1-\sigma\tau_{1}=\varphi_{1}\tau_{1}$ on each element,
\begin{equation}\label{eq:IandIV}
\begin{aligned}
  -\bigl(\textup{[I]}+\textup{[IV]}\bigr)
  &= \sigma\bigl(\tau_1\bv,\alpha\nabla p\bigr)_h
  + \bigl(p,\nabla\cdot(\alpha\bv)\bigr) \\
  &= \bigl(\varphi_1\tau_1\,p,\nabla\cdot(\alpha\bv)\bigr)_h
  + \sigma\sum_{b}\int_{\Gamma^b}\jump{\tau_1}\,\alpha\,(\bn\cdot\bv)\,p\dd\Gamma,
\end{aligned}
\end{equation}
where the sum extends over the interior faces $\Gamma^b$ of the partition: the elemental
boundary contributions have been assembled facewise using the continuity of
$\alpha,\bv,p$ across interior faces ($\bn$ is a fixed unit normal on each face and
$\jump{\tau_1}$ the corresponding jump), while the faces lying on $\Gamma$ vanish because
$\bv=\boldsymbol{0}$ there.

For the volumetric part of \eqref{eq:IandIV}, on each element
$\varphi_1\tau_1\le\varphi_1^{1/2}\tau_1^{1/2}$
(since $\varphi_1\tau_1\le1$, \cref{lem:parameters}\ref{P:basic}), and
$\varphi_1^{1/2}=c_1^{1/2}\alpha_K\tau_2^{1/2}/h$
(\cref{lem:parameters}\ref{P:27}), so that
\begin{equation}\label{eq:volpart}
\begin{aligned}
  \bigl|\bigl(\varphi_1\tau_1\,p,\nabla\cdot(\alpha\bv)\bigr)_h\bigr|
  &\le c_1^{1/2}\sum_K
  \tau_2^{1/2}\bigl\lVert\nabla\cdot(\alpha\bv)\bigr\rVert_K\;
  \alpha_K\,\tau_1^{1/2}h^{-1}\norm{p}_K \\
  &\le c_1^{1/2}\,\triplenorm{V}\,
  \bigl\lVert\alpha_K\tau_1^{1/2}h^{-1}p\bigr\rVert_h .
\end{aligned}
\end{equation}

For the jump part, fix an interior face $\Gamma^b$ shared by $K_1$ and $K_2$ and set
$\Omega^b=K_1\cup K_2$. By \cref{H:jump} the elemental quantities
$\varphi_1,\tau_1,\sigmatilde$ of $K_1$ and $K_2$ are uniformly comparable, so
\cref{lem:jump} yields, for any $i,j\in\{1,2\}$,
\begin{equation}\label{eq:jumpsplit}
  \sigma\bigl|\jump{\tau_1}\bigr|
  \le C\,\Bigl(\sigmatilde\tau_1\big|_{K_i}\Bigr)^{1/2}
        \Bigl(\sigmatilde\tau_1\big|_{K_j}\Bigr)^{1/2}
  \le C\,\sigmatilde^{1/2}\big|_{K_i}\;\tau_{1,K_j}^{1/2},
\end{equation}
the last step using $\sigmatilde\tau_1\le1$ (\cref{lem:parameters}\ref{P:sigmatilde})
on $K_j$ together with the comparability of $\tau_{1,K_i}$ and $\tau_{1,K_j}$.
Moreover $\alpha\le\alpha_{K_j}$ on $\Gamma^b\subset\overline{K}_j$. Bounding the face
integral by the $L^\infty$ norms over $\Omega^b$, expanding
$\norm{\bv}_{L^\infty(\Omega^b)}\norm{p}_{L^\infty(\Omega^b)}
\le\sum_{i,j\in\{1,2\}}\norm{\bv}_{L^\infty(K_i)}\norm{p}_{L^\infty(K_j)}$, and applying
the $L^\infty$ inverse estimate \eqref{eq:InverseEstimateFiniteOrderNorm} on each element (the diameters of
$K_1,K_2$ being comparable by \cref{H:mesh}), together with
$\operatorname{meas}(\Gamma^b)\le C h^{d-1}$,
\begin{equation}\label{eq:jumppart}
  \sigma\Bigl|\int_{\Gamma^b}\jump{\tau_1}\alpha(\bn\cdot\bv)p\dd\Gamma\Bigr|
  \le C\sum_{i,j\in\{1,2\}}
  \sigmatilde^{1/2}\big|_{K_i}\norm{\bv}_{K_i}\;
  \alpha_{K_j}\,\tau_{1,K_j}^{1/2}\,h^{-1}\norm{p}_{K_j}.
\end{equation}
Summing over the interior faces and applying the Cauchy--Schwarz inequality to the sums
(each element meets a uniformly bounded number of faces, \cref{H:mesh}),
\begin{equation}\label{eq:jumptotal}
  \sigma\sum_{b}\Bigl|\int_{\Gamma^b}\jump{\tau_1}\alpha(\bn\cdot\bv)p
  \dd\Gamma\Bigr|
  \le C\,\bigl\lVert\sigmatilde^{1/2}\bv\bigr\rVert_h\,
  \bigl\lVert\alpha_K\tau_1^{1/2}h^{-1}p\bigr\rVert_h
  \le C\,\triplenorm{V}\,
  \bigl\lVert\alpha_K\tau_1^{1/2}h^{-1}p\bigr\rVert_h .
\end{equation}
Combining \eqref{eq:volpart} and \eqref{eq:jumptotal},
\begin{equation}\label{eq:IIVbound}
  \bigl|\textup{[I]}+\textup{[IV]}\bigr|
  \le C\,\triplenorm{V}\,
  \bigl\lVert\alpha_K\tau_1^{1/2}h^{-1}p\bigr\rVert_h .
\end{equation}

\smallskip\noindent\emph{Step 6c: the terms \textup{[II]+[V]}.}
Again with $1-\sigma\tau_1=\varphi_1\tau_1$ elementwise, and using
$\varphi_1^{1/2}\tau_1\le\tau_1^{1/2}$ together with
\cref{lem:parameters}\ref{P:27},
\begin{align}
  \bigl|\textup{[II]}+\textup{[V]}\bigr|
  = \bigl|\bigl(\varphi_1\tau_1\,\bu,\alpha\X(V)\bigr)_h\bigr|
  &\le \sum_K \varphi_1^{1/2}\norm{\bu}_K\;
  \varphi_1^{1/2}\tau_1\bigl\lVert\alpha\X(V)\bigr\rVert_K
  \nonumber\\
  &\le c_1^{1/2}\sum_K \alpha_K\tau_2^{1/2}h^{-1}\norm{\bu}_K\;
  \tau_1^{1/2}\bigl\lVert\alpha\X(V)\bigr\rVert_K
  \nonumber\\
  &\le c_1^{1/2}\,
  \bigl\lVert\alpha_K\tau_2^{1/2}h^{-1}\bu\bigr\rVert_h\,\triplenorm{V}.
  \label{eq:IIVterm}
\end{align}

\smallskip\noindent\emph{Step 6d: the inter-element term \textup{[III]}.}
As in Step~6b, the elemental boundary contributions assemble facewise (the integrand
$\alpha(\bn\cdot\ba)(\bu\cdot\bv)$ is single-valued on interior faces and $\bu$ vanishes
on $\Gamma$):
\begin{equation*}
  \textup{[III]}
  = -\sigma\sum_b\int_{\Gamma^b}\jump{\tau_1}\,
  \alpha(\bn\cdot\ba)(\bu\cdot\bv)\dd\Gamma .
\end{equation*}
On $\Gamma^b$ we bound $\alpha\,|\bn\cdot\ba|\le\alpha_{K_j}|\ba|_{\infty,K_j}
\le\varphi_{1,K_j}h_{K_j}/c_2$ (\cref{lem:parameters}\ref{P:components}). Using
\eqref{eq:jumpsplit} in the symmetric form
$\sigma|\jump{\tau_1}|\le C(\sigmatilde\tau_1|_{K_i})^{1/2}
(\sigmatilde\tau_1|_{K_j})^{1/2}$, and
$\tau_{1,K_i}^{1/2}\tau_{1,K_j}^{1/2}\varphi_{1,K_j}
=\bigl(\tau_{1,K_i}^{1/2}\varphi_{1,K_j}^{1/2}\bigr)
\bigl(\tau_1\varphi_1|_{K_j}\bigr)^{1/2}\le C$ (comparability of $\tau_1$ across the face
and $\varphi_1\tau_1\le1$, \cref{lem:parameters}\ref{P:basic}), we obtain on each face, after applying the $L^\infty$ inverse
estimates and $\operatorname{meas}(\Gamma^b)\le Ch^{d-1}$ exactly as in Step~6b,
\begin{equation}\label{eq:IIIface}
  \sigma\Bigl|\int_{\Gamma^b}\jump{\tau_1}\alpha(\bn\cdot\ba)(\bu\cdot\bv)
  \dd\Gamma\Bigr|
  \le C\sum_{i,j\in\{1,2\}}
  \sigmatilde^{1/2}\big|_{K_i}\norm{\bu}_{K_i}\;
  \sigmatilde^{1/2}\big|_{K_j}\norm{\bv}_{K_j},
\end{equation}
whence, summing over faces and applying the Cauchy--Schwarz inequality as before,
\begin{equation}\label{eq:IIIbound}
  \bigl|\textup{[III]}\bigr|
  \le C\,\bigl\lVert\sigmatilde^{1/2}\bu\bigr\rVert_h
  \bigl\lVert\sigmatilde^{1/2}\bv\bigr\rVert_h
  \le C\,\triplenorm{U}\,\triplenorm{V}.
\end{equation}
Collecting \eqref{eq:T13conv}, \eqref{eq:IIVbound}, \eqref{eq:IIVterm} and
\eqref{eq:IIIbound} (and recalling that $R$ has been absorbed into $N$, so that
$T_{13}+T_2+T_4+T_{11}=T_{13}^{\mathrm{c}}+N$),
\begin{equation}\label{eq:groupstep}
  |T_{13}^{\mathrm{c}}| + |N|
  \le C\,\triplenorm{V}\Bigl(\triplenorm{U}
  + \bigl\lVert\alpha_K\tau_2^{1/2}h^{-1}\bu\bigr\rVert_h
  + \bigl\lVert\alpha_K\tau_1^{1/2}h^{-1}p\bigr\rVert_h\Bigr).
\end{equation}

\medskip\noindent\emph{Step 7: the penalty cross terms.}
By \cref{lem:parameters}\ref{P:eps}, on each element
$\varepsilon\tau_2^{1/2}=(\varepsilon\tau_2)^{1/2}\varepsilon^{1/2}
\le C_2^{1/2}\varepsilon^{1/2}$, hence
\begin{equation}\label{eq:crosspenalty}
  |T_{16}|
  \le \bigl\lVert\tau_2^{1/2}\nabla\cdot(\alpha\bv)\bigr\rVert_h\;
  \varepsilon\bigl\lVert\tau_2^{1/2}p\bigr\rVert_h
  \le C_2^{1/2}\,\triplenorm{V}\;\varepsilon^{1/2}\norm{p}
  \le \triplenorm{V}\,\triplenorm{U},
\end{equation}
and likewise $|T_{17}|\le\triplenorm{U}\,\triplenorm{V}$.

\medskip\noindent\emph{Step 8: assembly.}
Adding the estimates of Steps 1--7,
\begin{equation}\label{eq:assembly}
  \bigl\lvert\BASGS (\ba;U,V)\bigr\rvert
  \le C\,\triplenorm{V}\Bigl(\triplenorm{U}
  + \bigl\lVert\alpha_K\tau_2^{1/2}h^{-1}\bu\bigr\rVert_h
  + \bigl\lVert\alpha_K\tau_1^{1/2}h^{-1}p\bigr\rVert_h\Bigr).
\end{equation}

\medskip\noindent\emph{Step 9: absorption of the working norm.}
It remains to bound $\triplenorm{U}$ for \emph{discrete} $U=U_h$ by the bracketed norms.
We estimate each contribution to \eqref{eq:triplenorm} elementwise, using
\cref{lem:winv,lem:parameters}:
\begin{align}
  \nu^{1/2}\bigl\lVert\alpha^{1/2}\ViscProj\nabla\bu\bigr\rVert_K
  &\le \nu^{1/2}\frac{\Cinv}{h}\alpha_K^{1/2}\norm{\bu}_K
  \le \frac{\Cinv}{c_1^{1/2}}\,\varphi_1^{1/2}\norm{\bu}_K
  = \Cinv\,\alpha_K\,\tau_2^{1/2}h^{-1}\norm{\bu}_K,
  \label{eq:absorb1}\\
  \bigl\lVert\sigmatilde^{1/2}\bu\bigr\rVert_K
  &\le \varphi_1^{1/2}\norm{\bu}_K
  = c_1^{1/2}\,\alpha_K\,\tau_2^{1/2}h^{-1}\norm{\bu}_K,
  \label{eq:absorb2}\\
  \varepsilon^{1/2}\norm{p}_K
  &\le c_1^{1/2}\,\alpha_K\,\tau_1^{1/2}h^{-1}\norm{p}_K,
  \label{eq:absorb3}\\
  \tau_1^{1/2}\bigl\lVert\alpha\X(U)\bigr\rVert_K
  &\le \tau_1^{1/2}\alpha_K|\ba|_\infty\frac{\Cinv}{h}\norm{\bu}_K
   + \tau_1^{1/2}\alpha_K\frac{\Cinv}{h}\norm{p}_K \nonumber\\
  &\le \frac{\Cinv\,c_1^{1/2}}{c_2}\,\alpha_K\tau_2^{1/2}h^{-1}\norm{\bu}_K
   + \Cinv\,\alpha_K\,\tau_1^{1/2}h^{-1}\norm{p}_K,
  \label{eq:absorb4}\\
  \tau_2^{1/2}\bigl\lVert\nabla\cdot(\alpha\bu)\bigr\rVert_K
  &\le \Cinva\,\alpha_K\,\tau_2^{1/2}h^{-1}\norm{\bu}_K .
  \label{eq:absorb5}
\end{align}
Here \eqref{eq:absorb1} uses \eqref{eq:winv-grad} and
\ref{P:components}--\ref{P:27}; \eqref{eq:absorb2} uses
\ref{P:sigmatilde} and \ref{P:27}; \eqref{eq:absorb3} uses \ref{P:eps};
in \eqref{eq:absorb4} the velocity part uses \eqref{eq:winv-conv},
$\tau_1^{1/2}\alpha_K|\ba|_\infty/h \le \varphi_1\tau_1^{1/2}/c_2
\le\varphi_1^{1/2}/c_2$ and \ref{P:27}, while the pressure part uses
\eqref{eq:winv-gradp} directly; and \eqref{eq:absorb5} uses \eqref{eq:winv-divu}. Squaring
and summing over the elements,
\begin{equation}\label{eq:normconv}
  \triplenorm{U_h}
  \le C\Bigl(
  \bigl\lVert\alpha_K\tau_2^{1/2}h^{-1}\buh\bigr\rVert_h
  + \bigl\lVert\alpha_K\tau_1^{1/2}h^{-1}p_h\bigr\rVert_h\Bigr).
\end{equation}
Inserting \eqref{eq:normconv} in \eqref{eq:assembly} proves \eqref{eq:sharpcont}, and
$\alpha_K\le1$ \cref{H:porosity} gives \eqref{eq:continuity}.
\end{proof}

\begin{remark}[Where the porosity weight comes from]\label{rem:sharper}
The weight $\alpha_K$ of \eqref{eq:sharpcont}, which propagates to the convergence estimate
\eqref{eq:convergence}, survives every step of the proof for one reason worth recording: the
smallness condition \eqref{eq:UpperBoundOnEpsilon} on $\varepsilon$ carries the factor
$\alpha_K^2$, and that is exactly what the penalty contribution \eqref{eq:absorb3} needs in
order to carry the weight $\alpha_K$ as well. The same factor is the one demanded by the
stability analysis, so the two requirements coincide rather than compete.
\end{remark}

\begin{remark}[Where each hypothesis is used]\label{rem:hypotheses}
\Cref{H:data} (condition \eqref{eq:UpperBoundOnEpsilon}) enters Steps 3, 7 and 9;
\cref{H:projector} enters through the symmetric form of $T_1$ and through
\cref{lem:winv}, whose proof uses both the constant{\hyp}coefficient action of
$\ViscProj$ and the contraction $|\ViscProj\mathbf{T}|\le|\mathbf{T}|$ (Steps 1, 4, 5 and 9);
its Korn half \cref{eq:PiKorn} is \emph{not} used here, only in \cref{lem:definiteness};
\cref{H:porosity} enters through \cref{lem:winv} (Steps 4, 5 and 9) and through
$\alpha_K\le1$; \cref{H:advection} enters Step 6a; \cref{H:spaces} enters wherever
$\tau_1,\tau_2$ are moved in or out of an elemental integral, in particular in Steps 2 and 6a;
\cref{H:mesh} enters every inverse estimate and the facewise bookkeeping of Steps 6b and 6d;
\cref{H:jump} enters only \eqref{eq:jumpsplit}, i.e.\ the jump terms of Steps 6b and 6d, which
carry the factor $\sigma$ and hence vanish identically when $\sigma=0$. As
in~\cite{codina2001stabilized}, the jump condition and the $L^\infty$ inverse estimate are
therefore needed only when $\sigma>0$.
\end{remark}

\subsection{Convergence}\label{sec:convergence}
Stability and continuity are now combined in the standard way. The only new ingredient is that
the continuity estimate has to be re-run with the \emph{interpolation error} in its first slot,
where the inverse estimates of \cref{lem:winv} are no longer available and must be replaced by
interpolation estimates.

\subsubsection{Interpolation estimates}
Let $U=[\bu;p]$ denote the solution of the continuous linearized problem, whose regularity
\cref{H:regularity} we write, allowing for extra smoothness, as
\begin{equation}\label{eq:regularity}
  U\in\mathcal{W}_r \coloneqq
  \bigl(H^{m}(\Omega)\cap H^1_0(\Omega)\bigr)^d
  \times\bigl(H^{n}(\Omega)\cap\mathcal{Q}_0\bigr),
\end{equation}
with $m\ge k_u+1\ (\ge 2)$ and $n\ge k_p+1\ (\ge 2)$.
In particular
$\mathcal{L}_{\ba}U\in L^2(K)$ for every $K$, so that the consistency property
\eqref{eq:consistency} holds, and for $d\le3$ both $\bu$ and $p$ are continuous on
$\overline{\Omega}$. Let $\buhat\in\mathcal{V}_{h0}$ be the Lagrange interpolant of $\bu$ and
$I_h p$ that of $p$. For $\varepsilon>0$ we set $\widehat{p}_h\coloneqq I_h p$; when
$\mathcal{Q}_{h0}$ is the zero{\hyp}mean subspace---the case $\varepsilon=0$---the interpolant
must be mean{\hyp}corrected to
$\widehat{p}_h\coloneqq I_h p-\tfrac{1}{|\Omega|}\int_\Omega I_h p$ so that
$\widehat{p}_h\in\mathcal{Q}_{h0}$. In either case we set
$\widehat{U}_h\coloneqq[\buhat;\widehat{p}_h]\in\Xh$.

The mean shift is a \emph{global} constant, and therefore not controlled by any local
interpolation error; it is, however, invisible to the analysis exactly in the case where it is
applied. Indeed, the first{\hyp}slot pressure enters \eqref{eq:Bstab} only through $\nabla p$
(inside $\X$, \cref{eq:XG}) and through the two terms carrying $\varepsilon$, and it enters
\eqref{eq:triplenorm} only in the same two ways, so that for every constant $c$
\begin{equation}\label{eq:constshift}
  \BASGS \bigl(\ba;[\boldsymbol{0};c],V_h\bigr)=0
  \quad\text{and}\quad
  \triplenorm{[\boldsymbol{0};c]}=0
  \qquad (\varepsilon=0,\ \forall\, V_h\in\Xh).
\end{equation}
All the estimates below may therefore be carried out with the \emph{uncorrected} interpolant
$I_h p$, the correction being reinstated only where the membership
$\widehat{U}_h-U_h\in\Xh$ is invoked, in the proof of \cref{thm:convergence}. We accordingly
write the interpolation estimates for $I_h p$.

Let the elemental interpolation error
\cref{eq:InterpolationError} of the main text be evaluated on the two fields,
\begin{equation}\label{eq:Eint}
  \Eint{K}{\bu}\coloneqq h_K^{k_u+1}\norm{\bu}_{H^{k_u+1}(K)},
  \qquad
  \Eint{K}{p}\coloneqq h_K^{k_p+1}\norm{p}_{H^{k_p+1}(K)} .
\end{equation}
Standard interpolation theory~\cite{brenner2008mathematical,ernguermond} provides, on every $K$,
\begin{equation}\label{eq:interp}
\begin{aligned}
  \norm{\bu-\buhat}_{H^{m'}(K)} &\le C\,h_K^{-m'}\,\Eint{K}{\bu},
  \qquad m'\in\{0,1,2\},\\
  \norm{p-I_h p}_{H^{m'}(K)} &\le C\,h_K^{-m'}\,\Eint{K}{p},
  \qquad m'\in\{0,1\},
\end{aligned}
\end{equation}
and, since $k_u+1,\,k_p+1>d/2$ for $d\le3$, also the $L^\infty$ estimates
\begin{equation}\label{eq:interpinfty}
\begin{aligned}
  \norm{\bu-\buhat}_{L^\infty(K)} &\le C\,h_K^{-d/2}\,\Eint{K}{\bu},\\
  \norm{p-I_h p}_{L^\infty(K)} &\le C\,h_K^{-d/2}\,\Eint{K}{p},
\end{aligned}
\end{equation}
which are the analogues of Eqs.~(78)--(79) of~\cite{codina2001stabilized}. We abbreviate
$E\coloneqq[\beu;e_p]\coloneqq[\bu-\buhat;\,p-I_h p]$, interchangeable with $U-\widehat{U}_h$
throughout by \eqref{eq:constshift}, and define
\begin{equation}\label{eq:psih}
\begin{aligned}
  \Psi(h)
  &\coloneqq \Bigl(\sum_{K\in\Th}\alpha_K^2\,h_K^{-2}
  \bigl(\tau_{2,K}\,\Eint{K}{\bu}^2+\tau_{1,K}\,\Eint{K}{p}^2\bigr)\Bigr)^{1/2} \\
  &\;\le\;
  \overline{\Psi}(h) \coloneqq \sum_{K\in\Th}\alpha_K\,h_K^{-1}
  \Bigl(\tau_{2,K}^{1/2}\,\Eint{K}{\bu}+\tau_{1,K}^{1/2}\,\Eint{K}{p}\Bigr) \\
  &\;\le\;
  \sum_{K\in\Th}h_K^{-1}
  \Bigl(\tau_{2,K}^{1/2}\,\Eint{K}{\bu}+\tau_{1,K}^{1/2}\,\Eint{K}{p}\Bigr),
\end{aligned}
\end{equation}
the first inequality being the discrete $\ell^2\subset\ell^1$ bound (using
$a^2+b^2\le(a+b)^2$ elementwise). Here $\Psi(h)$ is the porosity{\hyp}weighted $\ell^2$
functional appearing on the right of the main{\hyp}text estimate \eqref{eq:ConvergenceResult}---we
drop its method subscript, as we do for the norm---and $\overline{\Psi}(h)$ is its coarser
$\ell^1$ majorant. The estimates below are therefore stated in
$\Psi(h)$: on a quasi{\hyp}uniform family the passage to $\ell^1$ can cost a mesh{\hyp}cardinality
factor $\sim h^{-d/2}$, so the optimal orders must be read off the $\ell^2$ functional.

\subsubsection{Continuity with respect to the interpolation error}

\begin{lemma}[Continuity for the interpolation error]\label{lem:continterp}
Let \cref{H:mesh,H:data,H:projector,H:porosity,H:advection,H:spaces} and the jump condition \cref{H:jump} (when $\sigma>0$) hold, and let $U\in\mathcal{W}_r$. Then
\begin{equation}\label{eq:continterp}
  \bigl|\BASGS \bigl(\ba;\,U-\widehat{U}_h,\,V_h\bigr)\bigr|
  \le C\,\Psi(h)\,\triplenorm{V_h}
  \qquad\forall\,V_h\in\Xh ,
\end{equation}
with $\Psi(h)$ the functional of \eqref{eq:psih}.
\end{lemma}

\begin{proof}
We rerun the proof of \cref{lem:continuity} with first argument $E=[\beu;e_p]$ and second
argument $V_h$, which remains discrete and is treated exactly as before. Two
observations make this possible.

\smallskip\noindent\emph{(1) The exact identities remain valid.}
The decomposition of Step~0 and the identities of Steps 2, 3 and 6a--6b only require of
the first argument that $\beu\in H^1_0(\Omega)^d$ with $\beu|_K\in H^2(K)^d$, and
$e_p\in H^1(\Omega)$, together with the continuity of $\beu$ and $e_p$ across interior
faces; all of this holds by \eqref{eq:regularity} (for $d\le 3$,
$H^2(K)\hookrightarrow C^0(\overline{K})$ and $\bu,p\in C^0(\overline{\Omega})$, while
$\buhat,\widehat{p}_h$ are continuous finite element functions). In particular this places $E$
within the scope of \cref{lem:globalid}---which was stated for $H^1_0(\Omega)^d\times
H^1(\Omega)$ precisely so as to cover the non{\hyp}discrete first argument---so that
\eqref{eq:skew}--\eqref{eq:globalibp} apply verbatim, and the elementwise identities
\eqref{eq:elemibp} and those preceding \eqref{eq:IandIV} hold with well-defined traces.

\smallskip\noindent\emph{(2) Inverse estimates on the first argument are replaced by
interpolation estimates.}
Inspection of the proof of \cref{lem:continuity} shows that the first argument enters the
estimates only through the following elemental quantities; next to each we record the
discrete estimate used there and its replacement, obtained from \eqref{eq:resolved},
\eqref{eq:interp}--\eqref{eq:interpinfty} and
$|\ViscProj\mathbf{T}|\le|\mathbf{T}|$:
\begin{align}
  \bigl\lVert\nabla\cdot(\alpha\ViscProj\nabla\beu)\bigr\rVert_K
  &\le \alpha_{\infty,K}\bigl\lVert\nabla\cdot(\ViscProj\nabla\beu)\bigr\rVert_K
   + \norm{\nabla\alpha}_{L^\infty(K)}\bigl\lVert\ViscProj\nabla\beu\bigr\rVert_K \\
  &\le C\,\alpha_K\,h^{-2}\,\Eint{K}{\bu},
  \label{eq:interpdivvisc}\\[2pt]
  \bigl\lVert\alpha^{1/2}\ViscProj\nabla\beu\bigr\rVert_K
  &\le \alpha_K^{1/2}\,|\beu|_{H^1(K)} \\
  &\le C\,\alpha_K^{1/2}\,h^{-1}\,\Eint{K}{\bu},
  \label{eq:interpgrad}\\
  \bigl\lVert\nabla\cdot(\alpha\beu)\bigr\rVert_K
  &\le C\,\alpha_K\,h^{-1}\,\Eint{K}{\bu}, \\
   \norm{\alpha\,\ba\cdot\nabla\beu}_K
   &\le C\,\alpha_K\,|\ba|_{\infty,K}\,h^{-1}\,\Eint{K}{\bu},
  \label{eq:interpfirst}\\
  \norm{\alpha\nabla e_p}_K
  &\le C\,\alpha_K\,h^{-1}\,\Eint{K}{p}, \\
   \norm{\beu}_K &\le C\,\Eint{K}{\bu},
  \qquad
   \norm{e_p}_K \le C\,\Eint{K}{p},
  \label{eq:interpzero}\\
  \norm{\beu}_{L^\infty(K)}
  &\le C\,h^{-d/2}\,\Eint{K}{\bu}, \\
  \norm{e_p}_{L^\infty(K)} &\le C\,h^{-d/2}\,\Eint{K}{p}.
  \label{eq:interpinftyE}
\end{align}
These have exactly the same elemental coefficients as the weighted inverse estimates of
\cref{lem:winv} applied in Steps 4, 5, 6 and 9, with $\norm{\buh}_K$ and $\norm{p_h}_K$
replaced by $\Eint{K}{\bu}$ and $\Eint{K}{p}$. Consequently:
\begin{itemize}[leftmargin=2em]
\item Step 9 (the elemental bounds \eqref{eq:absorb1}--\eqref{eq:absorb5}) carried out
with \eqref{eq:interpdivvisc}--\eqref{eq:interpinftyE} gives
\begin{equation}\label{eq:Enorm}
  \triplenorm{E}
  \le C\Bigl(\sum_{K}\alpha_K^2\,h_K^{-2}\bigl(\tau_{2,K}(\Eint{K}{\bu})^2
  + \tau_{1,K}(\Eint{K}{p})^2\bigr)\Bigr)^{1/2}
  = C\,\Psi(h),
\end{equation}
the middle expression being $\Psi(h)$ of \eqref{eq:psih} written out.
\item The bracketed norms in \eqref{eq:assembly} are broken mesh $\ell^2${\hyp}norms and are
bounded, by the $m'=0$ estimates in \eqref{eq:interpzero}, as
$\bigl\lVert\alpha_K\tau_2^{1/2}h^{-1}\beu\bigr\rVert_h
+\bigl\lVert\alpha_K\tau_1^{1/2}h^{-1}e_p\bigr\rVert_h\le C\,\Psi(h)$.
\item In the jump terms \eqref{eq:jumppart} and \eqref{eq:IIIface} the $L^\infty$ inverse
estimates applied to the \emph{first} argument are replaced by the $L^\infty$
interpolation estimates \eqref{eq:interpinftyE} (those applied to $\bvh$ are kept), with
identical powers of $h$.
\end{itemize}
Steps 1--8 then yield
$\bigl|\BASGS (\ba;E,V_h)\bigr|\le C\triplenorm{V_h}\bigl(\triplenorm{E}+C\Psi(h)\bigr)$, and
\eqref{eq:Enorm} concludes the proof.
\end{proof}

\subsubsection{Main theorem}

\begin{theorem}[Convergence]\label{thm:convergence}
Under the hypotheses of \cref{lem:coercivity,lem:continterp},
\begin{equation}\label{eq:convergence}
\begin{aligned}
  \triplenorm{U-U_h}
  &\le C\,\Psi(h) \\
  &= C\Bigl(\sum_{K\in\Th}\alpha_K^2\,h_K^{-2}
  \bigl(\tau_{2,K}\,\Eint{K}{\bu}^2+\tau_{1,K}\,\Eint{K}{p}^2\bigr)\Bigr)^{1/2}
  \;\le\; C\,\overline{\Psi}(h).
\end{aligned}
\end{equation}
\end{theorem}

\begin{proof}
Let $\Delta_h\coloneqq\widehat{U}_h-U_h\in\Xh$ (not to be confused with the total error
$E_h = U-U_h$ of \cref{th:Convergence}). By \cref{lem:coercivity},
consistency \eqref{eq:consistency} and \cref{lem:continterp},
\begin{equation*}
  C_{\mathrm{coer}}\,\triplenorm{\Delta_h}^2
  \le \BASGS (\ba;\Delta_h,\Delta_h)
  = \BASGS \bigl(\ba;\widehat{U}_h-U,\Delta_h\bigr)
  \le C\,\Psi(h)\,\triplenorm{\Delta_h},
\end{equation*}
so $\triplenorm{\Delta_h}\le C\Psi(h)$. Here $\Delta_h$ is a legitimate test function precisely
because $\widehat{U}_h$ was mean{\hyp}corrected, which is the only role the correction plays;
$\triplenorm{U-\widehat{U}_h}=\triplenorm{E}$ by \eqref{eq:constshift}, and \eqref{eq:Enorm}
bounds the latter by $C\,\Psi(h)$. The triangle inequality then gives
$\triplenorm{U-U_h}\le\triplenorm{U-\widehat{U}_h}+\triplenorm{\Delta_h}\le C\,\Psi(h)\le C\,\overline{\Psi}(h)$.
\end{proof}

\noindent This is the porosity{\hyp}weighted form of the main{\hyp}text bound
\eqref{eq:ConvergenceResult}, whose interpretation is discussed there. We only add the caveat
already noted after \eqref{eq:psih}: the two relaxations---to $\overline{\Psi}(h)$, and then to
the porosity{\hyp}free weight $\alpha_K\equiv1$---are not innocuous in the same way. The second
merely enlarges the constant, but the first can cost a mesh{\hyp}cardinality factor
$\sim h^{-d/2}$, so the optimal convergence orders must be read off $\Psi(h)$ and not off its
$\ell^1$ majorant.
